\documentclass[11pt]{ctexart}
\usepackage[a4paper,margin=1in]{geometry}
\usepackage{graphicx}
\usepackage{xcolor}
\usepackage{hyperref}
\definecolor{refgray}{gray}{0.45}
\title{\textbf{市场最新观点汇总}\\[0.4em]{\large AI结构性需求与宏观再定价并行，能源、消费与地产分化加剧}}
\author{KC桌面 自动生成}
\date{260625}
\begin{document}
\maketitle
\tableofcontents
\newpage
\section{一页摘要}
\begin{itemize}
  \item AI算力与存储需求进入结构性扩张期，但市场对盈利可持续性的定价分歧加大，美股上涨逻辑从估值扩张切换至盈利驱动。
  \item 中国消费与地产数据出现裂口，一线城市二手房挂牌量持续下降为价格企稳提供支撑，但总量需求仍受制于财富效应消退。
  \item 能源与大宗商品领域，炼油利润韧性优于原油价格，铜矿股隐含定价过度悲观，焦煤供给结构性重置正在发生。
  \item 全球供应链运价高企从事件冲击转向供给约束，中国变压器出口受益于交付速度优势，日本造船业订单锁定至2029年。
  \item 欧洲央行加息进程未结束，德国养老金改革引入强制入市机制可能改变长期资产配置格局。
\end{itemize}
\section{宏观与利率：ECB加息未止，德国养老金改革重塑长期资金流}
\textbf{欧洲央行加息周期远未结束，规则框架指向终端利率高于市场定价；德国养老金改革引入强制资本积累制，每年或带来超800亿欧元增量资金，长期影响全球资产配置。}

\begin{itemize}
  \item DB通过九种货币政策规则检验ECB 6月加息，基于核心通胀的规则中位数指向终端利率2.75\%，意味着至少还需加息两次。
  \item ECB 6月上调通胀预测后，规则建议额外加息25-50个基点，但基于整体通胀的泰勒规则过于激进（建议年底加至4\%），在能源供给冲击下应忽略。
  \item 德国养老金改革提案中最具结构性意义的条款是引入强制性资本积累制，从2028年起逐步将2\%毛工资收入投入资本市场，预计2031年每年新增300-350亿欧元。
  \item DB估计同期私人养老金改革可能带来每年约500亿欧元资金流入，两者合计德国每年将有超800亿欧元新增资金涌入资本市场。
  \item 德国改革参考瑞典模式，在现收现付制之外强制劳动者将部分收入投入资本市场，通过长期复利积累个人养老金资产。
  \item 改革还包括将退休年龄从67岁逐步提至67.5岁、取消63岁提前退休优惠、将自由职业者纳入法定体系等务实措施。
  \item ECB 6月加息被拉加德称为“稳健”，DB规则分析支持这一判断，但市场对加息终点存在分歧，部分投资者认为周期尾声，另一部分担忧通胀粘性。
  \item 2020-2022年期间规则曾建议激进紧缩但央行未执行，2024年降息周期实际利率跟随规则下沿，表明规则指引与相机抉择之间存在张力。
\end{itemize}
\begin{figure}[htbp]
\centering
\includegraphics[width=0.88\linewidth]{figures/fig_104.jpg}
\caption{Figure 2 - Figure 2: Path of policy rates (ESTR) prescribed by a suite of monetary policy rules produced using the ECB's June staff baseline projections}
\end{figure}
\begin{figure}[htbp]
\centering
\includegraphics[width=0.88\linewidth]{figures/fig_105.jpg}
\caption{Figure 3 - Figure 3: Oil and gas futures contracts vs the ECB's June alternative scenarios}
\end{figure}
\begin{figure}[htbp]
\centering
\includegraphics[width=0.88\linewidth]{figures/fig_107.jpg}
\caption{Figure 4 - Figure 6: Severe scenario – policy rule prescriptions}
\end{figure}
\begin{figure}[htbp]
\centering
\includegraphics[width=0.88\linewidth]{figures/fig_112.jpg}
\caption{Figure 9 - Figure 9: ESTR vs policy rule prescriptions (\%) during the 2021-22 pandemic and energy shock}
\end{figure}
\begin{figure}[htbp]
\centering
\includegraphics[width=0.88\linewidth]{figures/fig_113.jpg}
\caption{Figure 10 - Figure 10: ESTR vs policy rule prescriptions (\%) during the latest rate cutting cycle}
\end{figure}
{\small\color{refgray}\textbf{References:} [R004] DB：ECB加息不是终点，而是规则框架下的“稳健”起点; [R029] DB：德国养老金改革的最大变量不是延迟退休，而是强制入市}

\section{AI/算力：需求三年翻五倍，MLCC与CoWoS成为新瓶颈}
\textbf{中国AI算力需求2028年将达2025年5倍，token消耗量指数级爆发；硬件瓶颈从GPU向MLCC和先进封装扩散，存储芯片长协谈判是估值重估关键。}

\begin{itemize}
  \item HSBC估算中国智能算力需求将从2025年约0.7GW增长至2028年约5.4GW，年复合增长率74\%，支撑来自token消耗量真实数据——2026年3月日均达140万亿。
  \item 中国超大规模云服务商资本开支2026年将突破7000亿元人民币，接近翻倍，但与美国仍有量级差距，增长空间依然可观。
  \item MS报告指出MLCC正从被动元件变为AI服务器主动瓶颈，单台AI服务器MLCC用量和价值量提升，全球前五大供应商占87\%份额，供给刚性。
  \item MS预计2027年全球CoWoS总需求逼近270万片，NVIDIA占比降至45\%但绝对值仍翻倍，CPU成为先进封装新变量，联发科意外崛起。
  \item JPM认为存储芯片处于“higher for longer”上行周期，核心驱动力从供需错配转向长期协议谈判，预计2026年下半年迎来密集宣布期。
  \item 美光财报显示季度收入同比增长346\%，已签署16份“照付不议”战略客户协议，五年累计承诺收入约1000亿美元，覆盖约20\% DRAM出货量。
  \item GS指出美光财报为欧洲半导体设备商ASML、ASMI、BESI及AI基础设施提供商NBIS提供需求可见性改善信号，但不同公司受益逻辑存在本质差异。
  \item NOM报告揭示中国AI投资驱动力是国家资本主导的“资金矩阵”，大基金三期累计规模近7000亿元，超长期特别国债首次将AI列为支持领域。
\end{itemize}
\begin{figure}[htbp]
\centering
\includegraphics[width=0.88\linewidth]{figures/fig_167.jpg}
\caption{Exhibit 4 - \#\# Strong compute demand to drive accelerated growth; we expect China intelligent compute to grow at a CAGR of \$74\textbackslash{}\%\$ over 2025-28e, reaching 5GW in 2028e With more advanced model launches in 2026, token consumption has shown strong growth. Per National Data A}
\end{figure}
\begin{figure}[htbp]
\centering
\includegraphics[width=0.88\linewidth]{figures/fig_168.jpg}
\caption{Exhibit 4 - As GAI penetration is still in the early stages, we expect the strong compute demand to persist and expect China intelligent compute (@FP16, or 16-bit binary floating-point computer number format) to grow to over 12 ZFLOPS (10\textasciicircum{}3 EFLOPS) in 2028e, close to 8x v}
\end{figure}
\begin{figure}[htbp]
\centering
\includegraphics[width=0.88\linewidth]{figures/fig_170.jpg}
\caption{Exhibit 6 - Exhibit 6. China CSPs continue to raise their capex Exhibit 7. We expect China intelligent compute to grow to c12,354 EFLOPS (@FP16) in 2028e Exhibit 8. We expect China intelligent compute data centre demand to grow at a CAGR of 74\% over 2025-28e and to reach }
\end{figure}
\begin{figure}[htbp]
\centering
\includegraphics[width=0.88\linewidth]{figures/fig_171.jpg}
\caption{Exhibit 7 - Exhibit 7. We expect China intelligent compute to grow to c12,354 EFLOPS (@FP16) in 2028e Exhibit 8. We expect China intelligent compute data centre demand to grow at a CAGR of 74\% over 2025-28e and to reach c5.4GW in 2028e Based on 1) our above forecast for C}
\end{figure}
\begin{figure}[htbp]
\centering
\includegraphics[width=0.88\linewidth]{figures/fig_242.jpg}
\caption{Exhibit 1 - Exhibit 1: Global CoWoS demand breakdown: 2026e vs. 2027e Global CoWoS capacity demand by key customer Exhibit 2: Global CoWoS demand Y/Y growth profile}
\end{figure}
\begin{figure}[htbp]
\centering
\includegraphics[width=0.88\linewidth]{figures/fig_243.jpg}
\caption{Exhibit 3 - Exhibit 3: Global CoWoS demand breakdown with newly introduced 2027e numbers Exhibit 4: Global CoWoS capacity expansion by year end and by vendor Exhibit 5: Global CoWoS consumption, by customer NVIDIA Broadcom AMD Xilinx AWS/A}
\end{figure}
\begin{figure}[htbp]
\centering
\includegraphics[width=0.88\linewidth]{figures/fig_234.jpg}
\caption{Figure 1 - Figure 1: CSP capex and AI memory \% share Figure 2: Relevant TAM \% of CSP capex Figure 3: Memory vs CSP market cap trend}
\end{figure}
\begin{figure}[htbp]
\centering
\includegraphics[width=0.88\linewidth]{figures/fig_235.jpg}
\caption{Figure 1 - Figure 1: CSP capex and AI memory \% share Figure 2: Relevant TAM \% of CSP capex Figure 3: Memory vs CSP market cap trend Figure 4: Absolute memory TAM \% of CSP capex}
\end{figure}
{\small\color{refgray}\textbf{References:} [R012] HSBC：中国AI算力需求三年翻五倍，这两家公司最受益; [R013] 摩根斯坦利：MS：AI硬件下一个瓶颈不是GPU，是MLCC; [R016] JPM：存储芯片的“更高更久”周期，真正考验的是长协落地速度; [R017] 摩根斯坦利：MS：2027年CoWoS的蛋糕，NVIDIA切走一半，但真正的变量在CPU; [R019] GS：美光财报暴露了一个被低估的欧洲AI供应链信号; [R011] NOM：中国AI投资的真正引擎不是BAT，是“国家资本矩阵”}

\section{能源与大宗：炼油利润韧性优于原油，铜矿股定价过度悲观}
\textbf{炼油产品利润率结构性支撑强于原油价格，铜矿股隐含定价低于现货铜价，焦煤供给因山西矿难发生结构性重置，定价锚从中国CFR转向澳大利亚FOB。}

\begin{itemize}
  \item GS报告显示美伊临时和平协议后原油价格暴跌超10美元，但炼油产品利润率降幅明显更小，全球炼油产品利润率指数仍维持战前水平两倍。
  \item GS预计2026Q4美国柴油利润46美元/桶、欧洲31美元/桶，仍为2013-2019年均值的2-3倍，2027年虽回落但仍处高位。
  \item 炼油利润韧性源于库存低位、波斯湾炼厂130万桶/日意外停产修复缓慢、以及中国和印度炼厂新增产能延迟等结构性支撑。
  \item GS铜业周材料指出多数铜矿股当前股价隐含铜价显著低于LME现货，南方铜业隐含铜价约11600美元/吨，比现货低约15\%。
  \item 市场在用“供过于求”剧本给铜矿股定价，但产业层面真实叙事是结构性短缺加深，供给约束才刚刚开始显现。
  \item GS焦煤报告指出山西矿难导致1500-2000万吨年度产量损失，部分矿井可能永远无法恢复满产，中国对优质进口焦煤依赖从周期性变为结构性。
  \item GS预计澳煤价格可能测试260-280美元/吨，中国港口价已上涨约40美元/吨，澳煤价格首次与国内煤持平。
  \item 印度炼焦煤进口预计从8000万吨增至1.2亿吨以上，但现货参与度低，定价权有限。
\end{itemize}
\begin{figure}[htbp]
\centering
\includegraphics[width=0.88\linewidth]{figures/fig_155.jpg}
\caption{Exhibit 8 - Exhibit 1: Refined Products Margins Have Fallen Less Since the Interim Peace Deal Announcement Than Crude Prices, With Our Global Refined Product Margin Index Still Double Its Pre-War Level Exhibit 2: We Have Nudged Down Our Di}
\end{figure}
\begin{figure}[htbp]
\centering
\includegraphics[width=0.88\linewidth]{figures/fig_156.jpg}
\caption{Exhibit 1 - Exhibit 4: Already Low Storage, a Stretched Refining System, and Rapidly Recovering Demand Will Likely Keep Product Stocks at Low Levels Over the Next 12 Months}
\end{figure}
\begin{figure}[htbp]
\centering
\includegraphics[width=0.88\linewidth]{figures/fig_159.jpg}
\caption{Exhibit 4 - Exhibit 6: Nearly 1.3mb/d of Persian Gulf Refining Capacity Remains Under Unplanned Maintenance/Repair More Than Two Months After the Ceasefire Announcement}
\end{figure}
\begin{figure}[htbp]
\centering
\includegraphics[width=0.88\linewidth]{figures/fig_161.jpg}
\caption{Exhibit 5 - Exhibit 7: We Expect Diesel and Gasoline Margins to Stay Slightly Above Market Forwards in 2026H2 and Remain Above 2025 Levels Through 2027 \#\# Two-Sided Risks to Margins, With a Smaller Downside Than for Crude}
\end{figure}
\begin{figure}[htbp]
\centering
\includegraphics[width=0.88\linewidth]{figures/fig_114.jpg}
\caption{Exhibit 1 - Exhibit 3: Over the past 3-4 years, New entrants have led to lower Asian prices, with Australia remaining uncompetitive, to incentivise imports into China China HCC netback to Aus QLD vs. FOB price (spot)}
\end{figure}
\begin{figure}[htbp]
\centering
\includegraphics[width=0.88\linewidth]{figures/fig_117.jpg}
\caption{Exhibit 3 - Exhibit 5: Australia remains a small share of China imports (\textbackslash{}\textasciitilde{}7\% in 2025), with Aus exports to China down materially from pre-2020 Australia Coking Coal Export Volumes to China Exhibit 6: Mongolian and Russian coal has almost fu}
\end{figure}
\begin{figure}[htbp]
\centering
\includegraphics[width=0.88\linewidth]{figures/fig_121.jpg}
\caption{Exhibit 7 - Exhibit 9: India met coal imports (Mt annualised) Global Coking Coal Prices Exhibit 10: Australian HCC, PCI and Semi-Soft prices}
\end{figure}
{\small\color{refgray}\textbf{References:} [R010] GS：炼油利润的韧性比原油价格更值得关注; [R003] GS铜业周：铜矿股正在交易一个“不存在的过剩”; [R005] GS：山西煤矿事故将结构性重塑全球焦煤定价权}

\section{供应链与运价：高运价从事件冲击转向供给约束，中国变压器交付速度成优势}
\textbf{全球运价高企核心驱动力从地缘事件切换为运力供给刚性约束，中国变压器出口受益于交付速度优势，日本造船业订单锁定至2029年。}

\begin{itemize}
  \item Bernstein报告指出海洋运价较冲突前上涨超100\%，空运运价高出40\%-50\%，但推动运价持续高位的底层逻辑正从供给冲击转向运力刚性约束。
  \item 美国陆运市场因监管执行趋严导致可用卡车运力持续收紧，现货TL运价（不含燃油）同比上涨48\%，监管因素带来的运力退出是结构性的。
  \item GS变压器出口追踪显示5月中国变压器总出口增速从27\%飙升至72\%，非洲（+252\%）、中东（+152\%）、美洲（+113\%）是主要拉动力。
  \item 美国本土变压器供需缺口高达72\%，GS预计到2028年收窄至57\%，但中国供应商交付速度是本土企业四倍（6-9个月 vs 128周）。
  \item GS指出日本造船业订单积压量达2927万总吨，相当于未来三年半产能已被填满，产能扩张政策与防务需求可能将景气周期从周期性修复升级为结构性重估。
  \item 日本国土交通省启动产能扩张框架，企业需制定从2024年水平扩张50\%以上产能计划才有资格获得3500亿日元振兴基金。
  \item 全球造船市场受伊朗局势影响运费飙升，新船价格仍处历史高位，日本船厂虽满负荷运转但政策支持有望打开产能天花板。
\end{itemize}
\begin{figure}[htbp]
\centering
\includegraphics[width=0.88\linewidth]{figures/fig_009.jpg}
\caption{EXHIBIT 3 - EXHIBIT 3: Global trade dashboard GLOBAL TRADE EXHIBIT 4: World trade volume growth EXHIBIT 5: US PMI and trade growth EXHIBIT 6: Value of US inventories per sector (excl. automotive)}
\end{figure}
\begin{figure}[htbp]
\centering
\includegraphics[width=0.88\linewidth]{figures/fig_010.jpg}
\caption{EXHIBIT 4 - EXHIBIT 4: World trade volume growth EXHIBIT 5: US PMI and trade growth EXHIBIT 6: Value of US inventories per sector (excl. automotive) EXHIBIT 7: Real value of US inventories per sector (excl. automotive), adjusted for US p}
\end{figure}
\begin{figure}[htbp]
\centering
\includegraphics[width=0.88\linewidth]{figures/fig_199.jpg}
\caption{Exhibit 1 - Exhibit 2: China total transformer export value growth picked up to \$72\textbackslash{}\%\$ yoy in May (vs \$27\textbackslash{}\%\$ in April)}
\end{figure}
\begin{figure}[htbp]
\centering
\includegraphics[width=0.88\linewidth]{figures/fig_200.jpg}
\caption{Exhibit 1 - Exhibit 2: China total transformer export value growth picked up to \$72\textbackslash{}\%\$ yoy in May (vs \$27\textbackslash{}\%\$ in April) Exhibit 3: China transformer export value to the US grew 114x yoy on a low base, or +26\% mom (vs +95\% yoy in April)}
\end{figure}
\begin{figure}[htbp]
\centering
\includegraphics[width=0.88\linewidth]{figures/fig_203.jpg}
\caption{Exhibit 4 - Exhibit 4: Export transformer to the US market has seen a relatively volatile ASP in 10-220MVA due to volatile mix change... US PPI for Electric Power and Specialty Transformer vs CPI (Jan 2018 indexed to 100) US PPI for electric}
\end{figure}
\begin{figure}[htbp]
\centering
\includegraphics[width=0.88\linewidth]{figures/fig_251.jpg}
\caption{Exhibit 3 - Exhibit 5: Monthly domestic order intake and order backlog (gross tonnage)}
\end{figure}
\begin{figure}[htbp]
\centering
\includegraphics[width=0.88\linewidth]{figures/fig_253.jpg}
\caption{Exhibit 4 - Exhibit 6: Global order backlog for alternative fuel capable ships (number of vessels) and their share of total order backlog Exhibit 7: Quarterly orders at Japanese shipbuilders and related companies}
\end{figure}
{\small\color{refgray}\textbf{References:} [R002] Bernstein：供应链的“高运价”不是短期脉冲，而是结构性重定价; [R014] GS：美国变压器缺的不是产能，是中国速度; [R020] GS：日本造船业被低估的三个催化剂，不只是AI的替代品}

\section{地产与消费：中国楼市拐点在挂牌量，消费数据出现5月裂口}
\textbf{中国一线城市二手房挂牌量持续下降为价格企稳构筑底部，但消费数据受政策脉冲消退和财富效应走弱拖累，奢侈品进口二季度明显降温。}

\begin{itemize}
  \item JPM周度数据显示一线城市二手房挂牌量从3月峰值下降2.6\%，上海单周挂牌量加速下降1.1\%，供给收缩正在改变定价权。
  \item 端午假期扰动周度数据，但剔除假期影响后10城二手房成交同比仍为正，市场尚未形成自我强化的内生动力。
  \item MS中国消费活动Z-Score五月继续走低，零售总额同比增速自2022年以来首次转负至-0.6\%，以旧换新政策效果正在消退。
  \item MS指出PPI转正主要驱动力是油价上涨而非国内需求复苏，非大宗商品部门利润率未改善，企业仍未重新获得定价权。
  \item DB数据显示中国奢侈品进口4-5月同比下滑7.2\%，一季度为正增长5.0\%，手袋进口从+5.1\%骤降至-11.2\%为最弱品类。
  \item 奢侈品进口走弱背后是房价持续承压（同比仍跌4.8\%）和股市财富效应快速消退（回报率从+21\%降至+14.8\%）。
  \item GS奶粉行业报告显示线下降幅边际收窄但量价结构尚未触底，伊利逆势转正（线下同比+1.9\%），飞鹤仍在追赶市场。
  \item 线上奶粉渠道5月降幅扩大至19\%，京东渠道同比下滑26\%，线上库存压力仍在释放中。
\end{itemize}
\begin{figure}[htbp]
\centering
\includegraphics[width=0.88\linewidth]{figures/fig_212.jpg}
\caption{Figure 1 - Figure 1: Iceberg 10-city real time secondary daily sales 冰山指数实时二手每日成交 (十大城市) Note: The Y/Y decline in the week ending 21 June 2026 is due to the Dragon Boat Festival.}
\end{figure}
\begin{figure}[htbp]
\centering
\includegraphics[width=0.88\linewidth]{figures/fig_213.jpg}
\caption{Figure 2 - Figure 2: Iceberg tier-1 cities secondary listing volume 冰山指数二手挂牌量 (一线城市) \#\# 2. Mainland China – leading indicators from Centraline Figure 3: Centraline secondary asking price index vs. NBS secondary home price index M/M in tier}
\end{figure}
\begin{figure}[htbp]
\centering
\includegraphics[width=0.88\linewidth]{figures/fig_218.jpg}
\caption{Figure 7 - Figure 7: 12-city daily secondary sales registrations (二手网签) Note: The decline in the week ending 21 June 2026 is due to the Dragon Boat Festival. Figure 8: 12-city secondary sales registrations (二手网签) 7-day moving average Secon}
\end{figure}
\begin{figure}[htbp]
\centering
\includegraphics[width=0.88\linewidth]{figures/fig_279.jpg}
\caption{Exhibit 1 - Exhibit 1: China Five-factor Consumer Activity Z-Score vs. MSCI China YoY change – Consumer activity declined further in May MS ASIA LIMITED+ Laura Wang Equity Strategist}
\end{figure}
\begin{figure}[htbp]
\centering
\includegraphics[width=0.88\linewidth]{figures/fig_258.jpg}
\caption{Exhibit 1 - Exhibit 1: Non-commodity sector profit margins have stabilized but remain relatively weak... Exhibit 2: ...leading to subdued trends in wage growth Exhibit 3: As the effects of the consumption trade-in program have faded, this}
\end{figure}
\begin{figure}[htbp]
\centering
\includegraphics[width=0.88\linewidth]{figures/fig_260.jpg}
\caption{Exhibit 2 - Exhibit 2: ...leading to subdued trends in wage growth Exhibit 3: As the effects of the consumption trade-in program have faded, this has resulted in an outright contraction in headline retail sales Exhibit 4: Hotel RevPAR indi}
\end{figure}
\begin{figure}[htbp]
\centering
\includegraphics[width=0.88\linewidth]{figures/fig_280.jpg}
\caption{Figure 2 - Figure 2: Limited improvement in lux imports YoY China monthly luxury import YoY Note: Mechanical watch import data as per NBS, Marh FHS data to be released on the 21st of April Figure 3: 1Q YoY improvement momentum did not carr}
\end{figure}
\begin{figure}[htbp]
\centering
\includegraphics[width=0.88\linewidth]{figures/fig_281.jpg}
\caption{Figure 3 - Figure 3: 1Q YoY improvement momentum did not carry into 2Q China quarterly luxury import (volume, YoY)}
\end{figure}
{\small\color{refgray}\textbf{References:} [R015] JPM：中国楼市真正的拐点，不在成交量，在挂牌量; [R026] 摩根斯坦利：中国消费的“5月裂口”比表面数据更值得警惕; [R023] 摩根斯坦利：MS：中国通胀的真相不在PPI，而在工资条; [R027] DB：中国奢侈品进口的“五月寒”比数据本身更值得警惕; [R028] GS：奶粉行业最坏的时候正在过去，但赢家已非旧面孔}

\section{区域市场：人民币中间价模型释放偏强信号，日本造船与热泵迎来结构性机会}
\textbf{人民币中间价模型预测值下调152个基点指向偏强方向，但政策通过逆周期因子平滑节奏；日本造船业订单锁定至2029年，英国热泵市场渗透率不足1\%但制冷需求爆发。}

\begin{itemize}
  \item NOM人民币中间价模型最新预测值为6.8043，较前次官方预测值6.8195低152个基点，计入逆周期因子后预测值6.8071，较前次定盘价低124个基点。
  \item 模型误差近期在收窄，说明模型对短期走势捕捉能力改善，但政策当局通过逆周期因子平滑节奏，非急转弯。
  \item 下半年关键事件密集：7月底政治局会议、11月APEC、12月中央经济工作会议、年底习主席访美，汇率弹性可能加大。
  \item GS日本造船业报告指出订单积压量达2927万总吨，相当于未来三年半产能已被填满，2026年前五个月累计订单量同比仅微降2\%。
  \item 日本造船业两大催化剂：产能扩张政策（3500亿日元基金）和防务需求（美国2027财年可能拨款18.5亿美元在日韩船厂造军舰）。
  \item GS热泵报告揭示全球变暖正在推动空气-空气热泵在英国加速渗透，2025年英国热泵销量约12.5万台同比增长27\%，渗透率仍不足1\%。
  \item 英国热泵市场政策补贴加码，空气源/地源热泵补贴7500英镑，但中国厂商通过低价+D2C模式渗透，欧洲本土OEM通过渠道护城河应对。
\end{itemize}
\begin{figure}[htbp]
\centering
\includegraphics[width=0.88\linewidth]{figures/fig_251.jpg}
\caption{Exhibit 3 - Exhibit 5: Monthly domestic order intake and order backlog (gross tonnage)}
\end{figure}
\begin{figure}[htbp]
\centering
\includegraphics[width=0.88\linewidth]{figures/fig_253.jpg}
\caption{Exhibit 4 - Exhibit 6: Global order backlog for alternative fuel capable ships (number of vessels) and their share of total order backlog Exhibit 7: Quarterly orders at Japanese shipbuilders and related companies}
\end{figure}
\begin{figure}[htbp]
\centering
\includegraphics[width=0.88\linewidth]{figures/fig_255.jpg}
\caption{Exhibit 6 - Exhibit 8: Quarterly operating profits for Japanese shipbuilding-related companies Furuno Electric's FY-end is February}
\end{figure}
{\small\color{refgray}\textbf{References:} [R024] NOM：人民币中间价模型已释放一个明确信号; [R020] GS：日本造船业被低估的三个催化剂，不只是AI的替代品; [R021] GS：热泵在英国卖的不是取暖，是空调}

\section{风险偏好：美股上涨逻辑从估值驱动切换至盈利驱动，AI牛市未能推升美元}
\textbf{标普500目标价上调至7800但估值假设下调，盈利增长成为核心驱动力；AI牛市未能同步推升美元，因美股跑赢含金量被高估且全球风险偏好回升稀释美元。}

\begin{itemize}
  \item BARC将标普500 2026年底目标价从7650上调至7800，同时将估值倍数从约24倍下调至23.1倍，EPS预期从321美元大幅上调至337美元。
  \item BARC认为推动美股上涨的发动机从流动性驱动、估值扩张转向盈利驱动、基本面支撑，一季度财报季科技板块EPS增速创2010年以来最佳。
  \item BARC指出AI资本开支峰值被推后，供给约束与需求曲线变陡并存，但超大规模企业自有现金流与资本开支缺口在扩大，2028年可能更紧张。
  \item GS报告揭示AI牛市未能推升美元的核心原因：美股跑赢更多是对发达市场而非新兴市场，盈利增长可持续性存疑，市场宽度太窄。
  \item GS指出全球股市（尤其亚洲科技重镇）也受益于AI行情，当全球风险偏好一起回升时美元通常承压，美股例外主义反而让美元更弱。
  \item BARC认为如果未来美股回调而全球其他市场相对抗跌，才是美元真正走弱的催化剂，目前AI行情对美元支撑有限。
  \item BARC将金融和医疗保健下调至中性，认为银行盈利上行空间有限，医疗保健估值已合理。
\end{itemize}
\begin{figure}[htbp]
\centering
\includegraphics[width=0.88\linewidth]{figures/fig_134.jpg}
\caption{FIGURE 2 - The earnings growth engine is firing on all cylinders. Activity data remain constructive: industrial production is on the rise, ISM manufacturing PMI is finally back in expansionary territory, and both durable goods and S\&P Global manufacturing output point to}
\end{figure}
\begin{figure}[htbp]
\centering
\includegraphics[width=0.88\linewidth]{figures/fig_135.jpg}
\caption{FIGURE 2 - FIGURE 2. The yield-equity correlation is testing the lower bound of its historical range at current levels of nominal 10Y FIGURE 3. The prospect of resumed hikes has not typically derailed equities ahead of the event S\&P 500 Performance Around Fed Reversal of}
\end{figure}
\begin{figure}[htbp]
\centering
\includegraphics[width=0.88\linewidth]{figures/fig_140.jpg}
\caption{FIGURE 8 - \#\# Raise FY26 EPS estimate to \textbackslash{}\$337 from \textbackslash{}\$321 We raise our FY26 S\&P 500 EPS estimate to \textbackslash{}\$337 from \textbackslash{}\$321, modestly below the Street's \textbackslash{}\$341 and implying 20.8\% Y/Y growth from \textbackslash{}\$279 in FY25. Since our last update, three developments argue for a higher EPS esti}
\end{figure}
\begin{figure}[htbp]
\centering
\includegraphics[width=0.88\linewidth]{figures/fig_142.jpg}
\caption{FIGURE 9 - Data as of 10 June 2026. FIGURE 9. We estimate \textbackslash{}\$337 in FY26 EPS (+21\% Y/Y) vs. Street at \textbackslash{}\$341 Data as of 16 Jun 2026 FIGURE 10. FY26 base, bull \& bear case EPS scenarios Data as of 10 Jun 2026 Looking ahead, we introduce a preliminary FY27 EPS estimate of \textbackslash{}\$}
\end{figure}
\begin{figure}[htbp]
\centering
\includegraphics[width=0.88\linewidth]{figures/fig_148.jpg}
\caption{FIGURE 17 - Our SOTP valuation framework assigns Big Tech a lower baseline multiple than in March (26x vs. 27.5x). We continue to believe the group offers durable earnings growth, but we trim our valuation assumptions to account for uncertainties around the scale, funding}
\end{figure}
\begin{figure}[htbp]
\centering
\includegraphics[width=0.88\linewidth]{figures/fig_149.jpg}
\caption{FIGURE 17 - Applying the blended multiple takes our YE26 S\&P 500 target up to 7800 from 7650, 23x our \textbackslash{}\$337 FY26 EPS estimate. Our upwardly revised EPS estimate does the heavy lifting, as our valuation assumptions are reduced modestly from our last update. Our base case r}
\end{figure}
\begin{figure}[htbp]
\centering
\includegraphics[width=0.88\linewidth]{figures/fig_001.jpg}
\caption{Exhibit 1 - Exhibit 2: US outperformance has been less pronounced versus EM equities than DM equities}
\end{figure}
\begin{figure}[htbp]
\centering
\includegraphics[width=0.88\linewidth]{figures/fig_004.jpg}
\caption{Exhibit 4 - Exhibit 5: US forward earnings growth expectations alone explain much of the residual between Dollar performance and relative equity performance Univariate regression of the USD\textbackslash{}\textasciitilde{}log(S\&P 500 / MSCI World ex-US) residual on US and}
\end{figure}
{\small\color{refgray}\textbf{References:} [R007] BARC：标普500目标价7800，但真正支撑上涨的不是估值，是盈利; [R001] GS：AI牛市没能让美元走强，问题出在“质量”而非“数量”}

\section{新能源车与电池：订单分化加速，电池行业进入三国分治阶段}
\textbf{中国新能源车订单出现冰火两重天，比亚迪单周破10万辆彰显规模效应；全球电池行业中国、韩国、日本走向不同战略路径，机器人电池是下一个价值高地。}

\begin{itemize}
  \item DB周度订单数据显示6月第三周比亚迪新订单10.68万辆环比暴增134\%，而蔚来环比-26\%、华为鸿蒙智行环比-36\%出现两位数下滑。
  \item 比亚迪订单暴增来自王朝+海洋双网体系全面开花，规模效应形成正向循环：更多订单意味着更强供应链议价权。
  \item 小米汽车环比+44\%至7200台，零跑环比+22\%至1.41万，理想环比+8\%但同比-34\%，特斯拉环比-4\%但同比+5\%。
  \item NOM报告指出全球电池需求仍在增长但结构质变，储能电池将以17\%年复合增速成为增长最快板块，2030年达926GWh。
  \item 中国企业凭借LFP电池和一体化供应链拿下全球78\%电动汽车电池和80\%储能电池市场份额，统治力在关税壁垒下仍在加固。
  \item 韩国企业转向高价值+本地化赛道，日本企业瞄准AI数据中心服务器电池，目标到FY29实现1000亿日元收入。
  \item NOM特别指出机器人电池单价是电动车电池3倍，预计2030年市场规模20-40亿美元，早期用高镍电池但LFP性能提升后可能切换。
\end{itemize}
\begin{figure}[htbp]
\centering
\includegraphics[width=0.88\linewidth]{figures/fig_271.jpg}
\caption{Figure 2 - Figure 2: Weekly new orders trend of key automakers (Li, NIO, XPeng, Leap, Xiaomi, Tesla) Figure 3: Weekly BYD, Geely, HIMA (mainly AITO) new orders trend Figure 4: Weekly HIMA (mainly AITO) new orders trend}
\end{figure}
\begin{figure}[htbp]
\centering
\includegraphics[width=0.88\linewidth]{figures/fig_272.jpg}
\caption{Figure 2 - Figure 2: Weekly new orders trend of key automakers (Li, NIO, XPeng, Leap, Xiaomi, Tesla) Figure 3: Weekly BYD, Geely, HIMA (mainly AITO) new orders trend Figure 4: Weekly HIMA (mainly AITO) new orders trend Figure 5: Weekly}
\end{figure}
\begin{figure}[htbp]
\centering
\includegraphics[width=0.88\linewidth]{figures/fig_273.jpg}
\caption{Figure 3 - Figure 3: Weekly BYD, Geely, HIMA (mainly AITO) new orders trend Figure 4: Weekly HIMA (mainly AITO) new orders trend Figure 5: Weekly Li Auto new orders trend Figure 6: Weekly NIO group new orders trend}
\end{figure}
{\small\color{refgray}\textbf{References:} [R025] DB：中国新能源车订单的“冰火两重天”已到关键分水岭; [R009] NOM：电池行业已进入“三国分治”阶段，机器人电池是下一个价值高地}

\section{外汇管理与机构配置：从被动不管理转向主动总组合}
\textbf{机构外汇管理正从后台被动合规提升为总组合层面主动策略，不管理汇率风险本身就是最昂贵的风险决策。}

\begin{itemize}
  \item DB报告指出机构投资者正将汇率风险管理从碎片化管理转向总组合方法，与超过35家机构客户讨论后确认这一趋势。
  \item 三个因素推动转变：荷兰养老金改革等区域监管变化、全球资产配置中非本币资产比重持续上升、机构对总组合方法的兴趣从理论转向落地。
  \item DB指出选择不对冲本身就是风险决策，若被动做出（因没人负责或太复杂），对组合结果影响往往不对称。
  \item 亚太地区机构在讨论中对汇率管理最为积极，因其非本币资产敞口增长最快。
  \item 远期仍是主力对冲工具，但期权越来越受欢迎，特别是那些必须完全对冲的基金用期权增加灵活性、降低整体成本。
  \item 机构需关注隐形敞口，例如持有美国股票但收入来自欧洲，汇率风险可能被低估。
\end{itemize}
\begin{figure}[htbp]
\centering
\includegraphics[width=0.88\linewidth]{figures/fig_154.jpg}
\caption{Figure 1 - Figure 1: Currency risk – the underrated opportunity under the total portfolio approach Quick win: Centralised FX can improve risk control without requiring a full portfolio reorganisation. Heatmap summarising the themes that fe}
\end{figure}
{\small\color{refgray}\textbf{References:} [R008] DB：机构外汇管理正从“被动不管理”转向“主动总组合”}

\section{稀土磁材与电力设备：中国之外产能规划漂亮但产出存疑，变压器交付速度成差异化优势}
\textbf{美国稀土永磁产能规划到2028年理论上覆盖当前进口需求，但客户认证、金属合金供应和成本竞争力决定实际产出；中国变压器交付速度是本土企业四倍。}

\begin{itemize}
  \item MS报告显示美国钕铁硼磁材产能到2028年可能达约2.8万吨/年，远超当前8500吨/年进口量，但产能不等于产出。
  \item 专家指出从工厂设计到满产有经验企业需3年，初创公司可能接近6年，客户认证、金属/合金供应、良品率是核心变量。
  \item 西方钕铁硼磁材质量有望3年内追上中国，但成本差距短期难缩小，专家估算西方本土产品可能比中国高30-50\%。
  \item GS变压器报告显示中国供应商交付周期6-9个月，而美国本土企业需128周，产能扩张周期中国约1年而美国需2-3年。
  \item GS对思源电气和国电南瑞给予买入评级，对华明装备维持中性，强调速度优势正在改写全球电力设备竞争规则。
  \item 美国电力变压器需求与本土供应缺口从当前72\%收窄至2028年57\%，非传统供应商仍有明确进入空间。
\end{itemize}
\begin{figure}[htbp]
\centering
\includegraphics[width=0.88\linewidth]{figures/fig_199.jpg}
\caption{Exhibit 1 - Exhibit 2: China total transformer export value growth picked up to \$72\textbackslash{}\%\$ yoy in May (vs \$27\textbackslash{}\%\$ in April)}
\end{figure}
\begin{figure}[htbp]
\centering
\includegraphics[width=0.88\linewidth]{figures/fig_200.jpg}
\caption{Exhibit 1 - Exhibit 2: China total transformer export value growth picked up to \$72\textbackslash{}\%\$ yoy in May (vs \$27\textbackslash{}\%\$ in April) Exhibit 3: China transformer export value to the US grew 114x yoy on a low base, or +26\% mom (vs +95\% yoy in April)}
\end{figure}
\begin{figure}[htbp]
\centering
\includegraphics[width=0.88\linewidth]{figures/fig_203.jpg}
\caption{Exhibit 4 - Exhibit 4: Export transformer to the US market has seen a relatively volatile ASP in 10-220MVA due to volatile mix change... US PPI for Electric Power and Specialty Transformer vs CPI (Jan 2018 indexed to 100) US PPI for electric}
\end{figure}
{\small\color{refgray}\textbf{References:} [R018] 摩根斯坦利：MS：稀土磁材的“中国之外”叙事，产能不是最大变量; [R014] GS：美国变压器缺的不是产能，是中国速度}

\section{结语}
综合来看，市场正处于结构性分化加速的窗口：AI硬件瓶颈从GPU向被动元件和先进封装扩散，能源利润韧性优于价格，中国消费与地产在政策脉冲消退后回归内生动力检验。投资者需关注盈利可见性、供给约束强度和长协落地速度等关键变量，而非简单追逐宏观叙事。
\newpage
\section{报告覆盖清单}
\begin{itemize}
  \item [R001] 风险偏好：美股上涨逻辑从估值驱动切换至盈利驱动，AI牛市未能推升美元 - GS：AI牛市没能让美元走强，问题出在“质量”而非“数量”
  \item [R002] 供应链与运价：高运价从事件冲击转向供给约束，中国变压器交付速度成优势 - Bernstein：供应链的“高运价”不是短期脉冲，而是结构性重定价
  \item [R003] 能源与大宗：炼油利润韧性优于原油，铜矿股定价过度悲观 - GS铜业周：铜矿股正在交易一个“不存在的过剩”
  \item [R004] 宏观与利率：ECB加息未止，德国养老金改革重塑长期资金流 - DB：ECB加息不是终点，而是规则框架下的“稳健”起点
  \item [R005] 能源与大宗：炼油利润韧性优于原油，铜矿股定价过度悲观 - GS：山西煤矿事故将结构性重塑全球焦煤定价权
  \item [R006] 未被主题摘要引用 - 摩根斯坦利：MS：中国5月成品油需求暴跌，电动车不是“冲击”，是“替代完成”
  \item [R007] 风险偏好：美股上涨逻辑从估值驱动切换至盈利驱动，AI牛市未能推升美元 - BARC：标普500目标价7800，但真正支撑上涨的不是估值，是盈利
  \item [R008] 外汇管理与机构配置：从被动不管理转向主动总组合 - DB：机构外汇管理正从“被动不管理”转向“主动总组合”
  \item [R009] 新能源车与电池：订单分化加速，电池行业进入三国分治阶段 - NOM：电池行业已进入“三国分治”阶段，机器人电池是下一个价值高地
  \item [R010] 能源与大宗：炼油利润韧性优于原油，铜矿股定价过度悲观 - GS：炼油利润的韧性比原油价格更值得关注
  \item [R011] AI/算力：需求三年翻五倍，MLCC与CoWoS成为新瓶颈 - NOM：中国AI投资的真正引擎不是BAT，是“国家资本矩阵”
  \item [R012] AI/算力：需求三年翻五倍，MLCC与CoWoS成为新瓶颈 - HSBC：中国AI算力需求三年翻五倍，这两家公司最受益
  \item [R013] AI/算力：需求三年翻五倍，MLCC与CoWoS成为新瓶颈 - 摩根斯坦利：MS：AI硬件下一个瓶颈不是GPU，是MLCC
  \item [R014] 供应链与运价：高运价从事件冲击转向供给约束，中国变压器交付速度成优势 / 稀土磁材与电力设备：中国之外产能规划漂亮但产出存疑，变压器交付速度成差异化优势 - GS：美国变压器缺的不是产能，是中国速度
  \item [R015] 地产与消费：中国楼市拐点在挂牌量，消费数据出现5月裂口 - JPM：中国楼市真正的拐点，不在成交量，在挂牌量
  \item [R016] AI/算力：需求三年翻五倍，MLCC与CoWoS成为新瓶颈 - JPM：存储芯片的“更高更久”周期，真正考验的是长协落地速度
  \item [R017] AI/算力：需求三年翻五倍，MLCC与CoWoS成为新瓶颈 - 摩根斯坦利：MS：2027年CoWoS的蛋糕，NVIDIA切走一半，但真正的变量在CPU
  \item [R018] 稀土磁材与电力设备：中国之外产能规划漂亮但产出存疑，变压器交付速度成差异化优势 - 摩根斯坦利：MS：稀土磁材的“中国之外”叙事，产能不是最大变量
  \item [R019] AI/算力：需求三年翻五倍，MLCC与CoWoS成为新瓶颈 - GS：美光财报暴露了一个被低估的欧洲AI供应链信号
  \item [R020] 供应链与运价：高运价从事件冲击转向供给约束，中国变压器交付速度成优势 / 区域市场：人民币中间价模型释放偏强信号，日本造船与热泵迎来结构性机会 - GS：日本造船业被低估的三个催化剂，不只是AI的替代品
  \item [R021] 区域市场：人民币中间价模型释放偏强信号，日本造船与热泵迎来结构性机会 - GS：热泵在英国卖的不是取暖，是空调
  \item [R022] 未被主题摘要引用 - NOM：字节跳动的音乐App，正在用AI绕过版权战争
  \item [R023] 地产与消费：中国楼市拐点在挂牌量，消费数据出现5月裂口 - 摩根斯坦利：MS：中国通胀的真相不在PPI，而在工资条
  \item [R024] 区域市场：人民币中间价模型释放偏强信号，日本造船与热泵迎来结构性机会 - NOM：人民币中间价模型已释放一个明确信号
  \item [R025] 新能源车与电池：订单分化加速，电池行业进入三国分治阶段 - DB：中国新能源车订单的“冰火两重天”已到关键分水岭
  \item [R026] 地产与消费：中国楼市拐点在挂牌量，消费数据出现5月裂口 - 摩根斯坦利：中国消费的“5月裂口”比表面数据更值得警惕
  \item [R027] 地产与消费：中国楼市拐点在挂牌量，消费数据出现5月裂口 - DB：中国奢侈品进口的“五月寒”比数据本身更值得警惕
  \item [R028] 地产与消费：中国楼市拐点在挂牌量，消费数据出现5月裂口 - GS：奶粉行业最坏的时候正在过去，但赢家已非旧面孔
  \item [R029] 宏观与利率：ECB加息未止，德国养老金改革重塑长期资金流 - DB：德国养老金改革的最大变量不是延迟退休，而是强制入市
\end{itemize}
\section{逐篇报告摘录}
\subsection{[R001] GS：AI牛市没能让美元走强，问题出在“质量”而非“数量”}
{\small\color{refgray}归类：风险偏好：美股上涨逻辑从估值驱动切换至盈利驱动，AI牛市未能推升美元}

2026年，市场最反直觉的宏观画面之一正在展开：美股在AI热潮中持续跑赢全球，但美元却未能同步走强。按照传统框架，这几乎是一个定价异常——美股的“例外主义”理应转化为对美元资产的需求，进而推高美元汇率。然而，GS外汇策略团队的最新报告指出，这一关系的断裂并非偶然，而是揭示了本轮AI驱动的美股上涨在“质量”上的根本缺陷。

这份由Lexi Kanter主笔的报告，核心判断出人意料：今年AI热潮带来的美股跑赢，很可能高估了美国经济的真实需求与美元资产的吸引力。美元之所以没有跟随美股起飞，不是因为市场不买账，而是因为这轮上涨的“纯度”不够——它太集中于少数公司、太依赖于短期盈利预期、且在全球范围内并未形成绝对优势。

这不仅仅是一个外汇定价的技术问题。它直接关系到投资者如何理解当前资产价格背后的驱动力，以及下半年宏观交易的核心矛盾。

KC评论： GS这份报告的关键洞察在于，它把“美股涨但美元不涨”这个现象拆解成了三个可验证的机制。对于多数读者来说，这比单纯看汇率点位更有价值——它提供了一套判断框架，让你能判断未来什么信号会真正改变美元走势。完整报告里还包含11条正在交易的策略建议，以及每个机制的详细回归检验图表，值得深入阅读。

1. 美股跑赢的“含金量”被高估：对新兴市场的优势远不及发达市场

GS首先指出，当前衡量美股相对表现的常用指标——标普500相对于MSCI全球除美国指数的比值——存在一个结构性失真。这个指数主要覆盖发达市场，而今年美股真正的跑赢幅度，在面对新兴市场时其实要小得多。

数据显示，韩国和台湾这两个科技权重极高的亚洲市场，在本轮AI资本开支浪潮中受益显著，其股市涨幅甚至在某些阶段超过了美股。这意味着，如果只盯着美股对发达市场的优势，会高估美国资产的相对吸引力。

更重要的是，资金流动对汇率的影响在不同市场间存在显著差异。GS的回归分析表明，美股资金流入对美元兑新兴市场货币汇率的驱动作用，远强于对美元兑发达市场货币汇率的影响。因此，当美股跑赢主要集中于发达市场一侧时，对美元的整体支撑自然有限。

这对投资者的含义很清晰：不能简单用“美股涨所以美元应该涨”来下注。AI红利正在全球范围内扩散，美国并非唯一赢家，美元的定价需要更精细的对手方分析。

2. 盈利增长的“持久性”才是关键：短期爆发无法替代长期预期

第二个机制可能更接近问题的本质。GS发现，美元对美股相对表现的敏感度，并非取决于盈利增长的绝对速度，而是取决于市场对增长持续性的预期。

具体而言，当市场预期美国盈利增长将在未来放缓——即一年期盈利预期显著高于两年期预期时——美元往往会跑输相对股市表现所暗示的水平。这正是当前的市场状态：AI热潮推动了一年期盈利预期的急剧上调，但两年期共识预期已经指向减速。

GS的回归模型显示，“美国自身盈利增速预期”这一变量，能够解释美元与相对股市表现之间残差的大部分。换句话说，市场并非不认可美股的基本面改善，而是质疑这种改善是否可持续。

KC评论： 这是一个值得所有资产配置者深思的结论。GS在报告中明确指出，盈利预期的“斜率”——即近端与远端的差距——比盈利水平本身更能解释汇率定价。这意味着，如果下半年AI盈利预期出现任何向下修正，美元面临的下行风险可能比市场想象的要大。完整报告中有一张图（Exhibit 4）清晰地展示了这一关系的动态演变，是理解当前交易逻辑的核心。

第三个机制则直指本轮AI牛市的结构性隐患。GS的股票策略团队指出，当前美国股市的宽度已经压缩至近几十年来最窄的水平之一——少数几只科技巨头贡献了指数的绝大部分涨幅。

GS用其内部的市场宽度指标进行检验后发现，当市场上涨越集中、宽度越窄时，美元相对于股市表现的“跑输”幅度往往越大。原因不难理解：少数股票的上涨所吸引的资本流入，更多集中在

[... middle omitted ...]

无动力”的均衡。

报告没有完全回答的是：什么事件能打破这一均衡？GS暗示，一个“美股下跌且全球其他市场明确跑赢”的场景，可能是触发美元贬值的催化剂。但问题在于，在当前AI叙事主导的环境下，这样的场景如何出现？是AI盈利预期的大面积下修？还是地缘政治冲击导致资金从风险资产全面撤离？抑或是美联储被迫转向更加鸽派的立场？

这些问题的答案，将决定下半年的宏观交易方向。而GS的11条活跃交易建议中，有两条直接反映了这一判断框架——做多韩国和台湾股票、做空印度和东南亚股票的组合，本质上是在押注AI红利的区域扩散将继续，这与美元承压的逻辑一脉相承。

综合GS的三个机制，投资者可以建立起一个更精细的美元分析框架：不要再问“美股涨了多少”，而要问“美股是怎么涨的”。

需要关注三个核心维度：第一，美股跑赢的对手方是谁——是发达市场还是新兴市场？第二，盈利增长是短期爆发还是长期趋势？第三，市场上涨的宽度是否在改善？只有当这三个维度都指向“高质量”的上涨时，美元才可能获得真正的支撑。

\subsection{[R002] Bernstein：供应链的“高运价”不是短期脉冲，而是结构性重定价}
{\small\color{refgray}归类：供应链与运价：高运价从事件冲击转向供给约束，中国变压器交付速度成优势}

Bernstein：供应链的“高运价”不是短期脉冲，而是结构性重定价

全球供应链正处于一个罕见的“三高”叠加窗口：运价高、需求韧、运力紧。这轮运价上涨并非单纯的地缘事件脉冲，而是正在重塑海运、空运、陆运的定价基准。Bernstein最新发布的《供应链脉搏报告》给出了一个与市场主流叙事略有不同的判断：高运价并未压制需求，消费者的购买行为正在适应新的运费成本，而供给侧的约束正在从临时性转向结构性。这份报告的核心洞察是，当前供应链的盈利环境对货运代理和运输公司而言是近期的“舒适区”，但不同环节、不同公司的受益程度将出现显著分化。

1. 这轮运价上涨的核心驱动力已从“事件冲击”切换为“供给约束”

报告明确指出，当前海洋运价较冲突前水平已上涨超过100\%，空运运价也高出40\%-50\%。但真正值得关注的不是涨幅本身，而是推动运价持续高位的底层逻辑正在发生切换。

最初的红海危机和苏伊士运河绕行是典型的供给冲击，但经过几个季度的消化，市场已经部分适应了更长的航线。当前运价持续高位的核心原因，是运力供给的刚性约束：海洋运输面临新船交付高峰前的运力错配，空运则受限于海湾地区承运人的运力瓶颈和燃油附加费，而美国陆运市场则因监管执行趋严（语言、居住地、运营范围、工时、培训等要求）导致可用的卡车运力持续收紧。

这意味着，即使地缘局势出现缓和，运价也不会迅速回到冲突前水平。供给侧的约束具有更强的粘性，尤其是在美国陆运市场，监管因素带来的运力退出是结构性的，而非周期性的。

KC评论： 市场往往习惯于用“地缘冲突缓和=运价回落”的线性逻辑来定价运输股，但Bernstein提醒我们，当前运价中有相当一部分是监管和运力结构变化带来的“硬成本”。这部分成本的回落速度会比想象中慢得多。完整报告中对美国陆运市场的监管细节和运力退出测算值得细读，这直接决定了联邦快递、UPS等公司的盈利中枢。

报告中最反直觉的发现之一是：运价上涨并未导致货运量萎缩。4月海洋集装箱量同比增长4\%，处于马士基全年2\%-4\%增长指引的上沿；空运货量也维持在0\%-5\%的增长区间。整体收入在5月和6月同比增长约40\%。

这背后是终端消费需求的韧性。尽管存在通胀压力，且油价仍高于冲突前水平，但消费者并未显著削减购买。Bernstein将此归结为供应链复杂性的持续存在——企业为了保障供应安全，愿意支付更高的运费来维持库存和交货可靠性。

这意味着，当前的高运价并非“量价齐跌”的恶性循环，而是“量稳价升”的良性窗口。对于货运代理和物流公司而言，这是一个难得的盈利环境：收入端受益于运价上涨，成本端则可以通过规模效应和灵活的运力配置来部分对冲。

3. 美国陆运市场正经历“供给推动”的运价飙升，需求端出现微弱曙光

美国陆运市场是报告重点分析的部分。现货市场卡车运价（不含燃油）同比涨幅已从5月的34\%进一步扩大至6月的48\%，连续三个月加速。Bernstein的判断是，这主要是一个供给驱动的故事，但需求端也出现了值得关注的积极信号。

供给侧的收紧主要来自监管执行：语言要求、居住地限制、运营范围界定、工时规定和培训标准等，正在系统性地淘汰不合规的承运人。美国最高法院对货运经纪人责任的裁决，虽然短期内不会立即冲击运价，但长期将进一步压缩运力。这应该会推动货运经纪行业的整合。

需求端则出现了“微光”：ISM制造业采购经理人指数连续多月处于扩张区间；托运人对制造业活动的信心有所改善（可能受到数据中心建设等工业活动的提振）；美国国内多式联运的年初至今累计货量在5月首次转为同比正增长（约+1\%）。

KC评论： 美国陆运市场正在经历一个“供给出清+需求弱复苏”的组合。对于投资者而言，关键在于区分哪些公司能在运价上涨中真正受益，哪些只是被动跟随。Bernstein对联邦快递和UPS给出

[... middle omitted ...]

士基是“跑输大盘”。

DSV被列为欧洲物流板块的首选。核心逻辑在于市场尚未充分定价其收购DB Schenker后的协同效应。DSV在过去20年的五笔并购中展现了出色的协同兑现能力，市场有理由期待它在新的大额交易中复制这一成功。一旦杠杆率回归到约2倍，公司有望重启大规模股份回购（稳态下每年可回购约240亿丹麦克朗），并开启新一轮并购。报告预计，整合完成后（2028年），DSV的每股收益将超过100丹麦克朗。

马士基则被明确看空。报告认为其核心业务——集装箱航运——面临挑战。红海危机虽然暂时阻止了2023年运价的崩溃，但这种“保护”正在消退。中东一旦达成持久和平协议，亚洲与欧洲之间的苏伊士航线恢复，将给运价带来新的压力。中期来看，行业订单簿上的新船数量庞大，在出现大规模拆船之前，运价环境将持续承压。更关键的是，Bernstein认为马士基的“整合商”战略存在缺陷，七年来未能兑现承诺的收入协同效应。

5. 报告尚未完全回答的关键问题：需求端的“微光”能否演变为“曙光”

\subsection{[R003] GS铜业周：铜矿股正在交易一个“不存在的过剩”}
{\small\color{refgray}归类：能源与大宗：炼油利润韧性优于原油，铜矿股定价过度悲观}

铜价横盘，库存高企，关税阴云未散——这是当前市场对铜板块最直观的感受。但GS刚刚发布的年度铜业周前瞻材料，给出了一个与市场情绪截然相反的判断：铜矿股的定价中隐含了太多悲观，而供给侧的真正约束才刚刚开始显现。

这不是一篇简单的会议预告。GS在材料中逐一拆解了全球主要铜矿上市公司的估值含义、增长路径和风险定价，得出的结论是：大多数铜矿股当前交易价格所隐含的铜价，已经低于现货水平。换句话说，市场在用“供过于求”的剧本给铜矿股定价，但产业层面的真实叙事却是“结构性短缺正在加深”。

这份材料的核心价值，不在于GS对铜价本身的预测，而在于它提供了一个极其清晰的框架：把每家公司股价倒推回其隐含的铜价假设，然后与现货铜价对比，就能直观判断市场在定价什么、又在忽略什么。

KC评论： GS这个“隐含铜价”的分析框架，是理解铜矿股当前定价锚点的最直接工具。如果一家公司股价隐含的铜价远低于现货，说明市场在用“衰退+过剩”的剧本定价；如果隐含铜价接近或高于现货，说明市场已经计入了供给紧张。完整报告中还包含了每家公司DCF模型的关键假设，如WACC、beta和终端增长率，这些参数对隐含铜价的敏感性极高，值得仔细对比。

1. 多数铜矿股定价已低于现货铜价，市场对供给过剩的担忧可能过度

GS覆盖的铜矿公司中，有多家当前股价所隐含的铜价显著低于LME现货铜价。以南方铜业为例，其隐含铜价约为每吨11,600美元，比现货低约15\%。Hudbay Minerals的隐含铜价约为每吨13,500美元，仅比现货低1\%，若计入Copper World项目则进一步降至11,700美元，折价幅度达14\%。First Quantum在计入Cobre Panama重启情景后，隐含铜价仅为每吨11,300美元，折价17\%。

这些数字意味着什么？市场在为这些公司定价时，已经假设了一个比当前更低的铜价环境。考虑到GS自身对铜市长期结构性短缺的判断，这种定价隐含的悲观情绪可能已经过度。

*对投资者的含义：** 如果铜价维持在当前水平甚至上行，这些被低估的标的将面临显著的估值修复空间。但需要注意，隐含铜价折价幅度最大的公司，往往也是运营风险或政治风险最高的公司——市场并非无缘无故给予折价。

GS在材料中反复提及一个主题：铜矿公司普遍面临成本上升压力，尤其是在能源和化学品领域。同时，矿石品位下降、监管变化和替代资本支出等结构性挑战正在持续压缩新项目的经济性。

以Freeport-McMoRan为例，GS预计其单位毛利将从2025年的每磅1.31美元提升至2027年的每磅2.90美元，增幅达121\%。这一增长并非来自铜价上涨，而是来自Grasberg矿的复产爬坡和浸出技术的推进。但市场对此仍持谨慎态度，导致FCX的估值低于其历史水平。

同样，Antofagasta的Centinela第二选矿厂（C2）项目已推进至50\%完成度，但市场尚未充分定价C2E等后续棕色地扩张选项。GS指出，若C2和C2E完全落地，Centinela可能成为行业规模最大、成本最低的运营资产之一。

*对投资者的含义：** 铜矿股的投资逻辑正在从“铜价涨跌”切换为“项目执行能力”。那些能够按时按预算交付扩张项目的公司，将在下一轮供给紧张中获得超额溢价。而成本超支或工期延误的公司，即使铜价上涨，也可能无法兑现利润。

KC评论： GS对Freeport的利润预测是一个很好的案例——它展示了当“产量增长”和“成本下降”同时发生时，利润弹性可以远高于铜价本身的涨幅。完整报告中还包含了每家公司 年的FCF yield预测，以及对应不同铜价情景的压力测试，这些数据是判断估值安全边际的关键。

3. Cobre Panama的重新启动是First Quantum的“期权价值”，但市场尚未定价

Fi

[... middle omitted ...]

Quantum现有资产的价值也已被过度折价。

*对投资者的含义：** First Quantum是一只高贝塔的“事件驱动型”铜矿股。核心问题是：Cobre Panama重启的概率到底有多大？GS没有给出明确概率，但材料中“encouraging developments”的措辞暗示，边际变化正在向积极方向移动。这一判断的验证，需要密切关注巴拿马政治环境和公司后续公告。

GS在材料中专门提到，拉丁美洲主要产铜国（秘鲁、智利）已度过一轮总统选举周期，巴西将在2025年下半年选举，阿根廷则在2027年。这在一定程度上降低了政治不确定性，但并不意味着监管风险消失。

以Southern Copper为例，其项目储备中包括Tia Maria（已获董事会批准）、El Arco、Los Chancas、Michiquillay等多个大型项目，但这些项目的推进速度、成本和社区关系，都高度依赖于东道国的监管环境。GS在风险部分明确列出了“墨西哥和秘鲁的政治与监管风险”作为下行因素。

\subsection{[R004] DB：ECB加息不是终点，而是规则框架下的“稳健”起点}
{\small\color{refgray}归类：宏观与利率：ECB加息未止，德国养老金改革重塑长期资金流}

ECB在6月会议上宣布加息25个基点，行长拉加德用“稳健”而非“强硬”来形容这一决定。市场对此的解读分化明显：一部分人认为这是本轮紧缩周期的尾声，另一部分人则担心通胀粘性将迫使ECB继续加码。DB最新发布的研报，通过一套严谨的政策规则框架，给出了一个清晰的判断：ECB的加息进程远未结束，但节奏和终点将由通胀结构而非总量决定。

这份研报的核心价值不在于预测具体加息几次，而在于它提供了一个可复用的分析框架——用ECB自己的经济预测来检验其政策路径的合理性。结论是：即使采用最温和的情景假设，ECB也需要至少再加息一次；而如果通胀从能源领域向核心领域扩散的假设成立，终端利率可能比市场当前定价高出25-50个基点。

这不仅仅是一份关于ECB的研报，它实际上在回答一个更根本的问题：在供给冲击主导的通胀环境中，中央银行应该如何校准政策？规则指引和相机抉择之间的张力，如何影响资产定价？

DB的分析起点非常直接：将ECB 6月员工预测数据输入九种不同的货币政策规则，看这些规则会给出什么样的利率路径。这九种规则包括泰勒规则、平衡法规则、前瞻性规则、惯性规则等，涵盖了对通胀指标（整体HICP vs 核心HICP）和产出缺口指标（失业率 vs 产出缺口）的不同选择。

结果是分层的。如果使用整体HICP作为通胀指标，泰勒型规则会建议ECB在年底前将利率大幅提升至4\%左右，这显然是一个不切实际的激进路径。但DB明确指出，在能源供给冲击的背景下，ECB完全可以忽略基于整体通胀的规则信号——因为货币政策无法直接影响能源价格本身，只能关注其溢出效应和第二轮效应。

更关键的是剔除了能源价格后的HICP。当使用这一指标时，规则给出的路径与当前市场定价的距离大幅收窄。中位数规则显示，到2026年底，政策利率应达到约2.75\%，这对应着两次额外加息。即便在ECB新引入的“温和情景”下，中位数规则也建议利率在年底达到略低于2.50\%的水平，对应一次额外加息。

KC评论： 这里的核心洞察是，ECB的加息决策不是基于对通胀的“恐惧”，而是基于一个可量化的规则框架。拉加德说“稳健”并非修辞，而是有数据支撑的。但读者需要注意，规则给出的只是“建议”，而非“剧本”。完整报告中详细展示了九种规则在三种情景下的具体路径，以及它们之间为何存在如此大的差异——这些细节对于理解ECB的实际决策边界至关重要。

DB还做了一项非常有价值的对比分析：比较3月和6月两次员工预测对规则路径的影响。结果显示，从3月到6月，通胀和增长预测的修正使得规则建议的政策利率路径平均提高了25-50个基点。

有趣的是，同期市场定价也恰好计入了一次额外的25个基点加息。这意味着，市场已经在某种程度上“消化”了ECB预测修正所隐含的政策收紧。但问题在于，这种消化是否充分？DB的规则框架表明，如果通胀从能源领域向核心领域扩散的假设持续成立，市场可能还需要进一步调整定价。

这一发现对于债券投资者尤其重要。它暗示当前收益率曲线中隐含的加息预期可能仍然偏低，尤其是在曲线的远端。如果ECB按照规则指引继续加息，长端利率可能面临上行压力。

KC评论： 这份对比分析的价值在于，它展示了政策规则如何帮助投资者“校准”对央行行为的预期。市场定价和规则指引之间的偏差，往往意味着交易机会或风险。完整报告中还包含了3月和6月预测修正的详细数据表，以及这些修正如何具体影响每一条规则的计算结果——这些细节是理解央行行为逻辑的关键。

DB非常清醒地指出，ECB从未机械地遵循任何一条规则。在 年疫情期间和俄乌冲突后的能源冲击中，规则曾建议ECB在2021年年中就开始加息，但实际利率直到2022年7月才开始上升。这种偏离并不一定是“政策错误”，因为央行需要考虑实时数据的滞后性、前瞻性判断以及金融稳定的考量。

但

[... middle omitted ...]

，但远端期货仍然更接近基线情景。

这意味着，ECB面临的是一个“时间错配”的挑战。短期内，能源价格下行有助于缓解通胀压力，从而降低进一步加息的必要性。但中期来看，如果能源期货价格重新上行，或者通胀从能源领域向核心领域扩散的势头持续，ECB将不得不重新评估其政策路径。

这里存在一个关键的未解问题：ECB的“温和情景”是否已经充分反映了美伊协议的影响？由于该协议的具体条款和执行细节尚未完全明确，ECB在6月预测中可能低估了其对能源价格的短期压制作用。如果能源价格在未来几个月持续低于预期，ECB可能会比规则建议的路径更早暂停加息。

KC评论： 这一点对于投资者来说可能是最大的“变量”。报告中的规则框架是基于ECB的官方预测，但官方预测本身存在滞后性。如果能源期货的实际走势持续低于ECB的温和情景假设，那么规则建议的加息路径可能就会“过时”。完整报告中包含了能源期货与ECB三种情景的对比图，以及不同情景下规则路径的详细拆解——这些图表是判断ECB政策“拐点”的重要工具。

\subsection{[R005] GS：山西煤矿事故将结构性重塑全球焦煤定价权}
{\small\color{refgray}归类：能源与大宗：炼油利润韧性优于原油，铜矿股定价过度悲观}

如果你只关心焦煤价格会涨到多少，这份GS报告给出的答案是每吨280美元。但真正值得你关注的，不是这个数字本身，而是它背后正在发生的三件事：中国正在失去高强度焦煤的自主供应能力，澳大利亚正在重新夺回对华定价权，而印度虽然需求增长最快，却始终无法成为价格的主导者。

这份基于新加坡铁矿石周和焦煤会议的调研报告，提供了一个罕见的全景视角。报告最核心的判断是：山西矿难导致的 万吨产量损失，不是一次性的供应冲击，而是一次结构性重置。部分矿井可能永远无法恢复到满产状态，中国对优质进口焦煤的依赖正在从周期性变为结构性。

这意味着，全球焦煤的定价锚正在从中国CFR价格，重新回到澳大利亚FOB指数。而这一转变，将影响从钢铁厂利润到矿业公司估值的整个链条。

1. 山西矿难不只损失了产量，还摧毁了中国对高强度焦煤的自给能力

GS团队在会议上得到的反馈显示，山西事故后仍有约 万吨产能处于停产状态。市场已经消化了 万吨的年度产量损失，但真正关键的问题是：这些产能还能回来吗？

与会交易商的判断并不乐观。部分矿井即使复产，产能利用率也可能只能恢复到之前的70-80\%，而过去有些矿井的实际产出甚至可以达到名义产能的110\%。监管问责的升级和地方官员的调查，正在让复产审批变得更加谨慎。

KC评论： 这里需要理解一个关键细节——损失的不仅是产量，更是特定品质的供应。山西事故中停产的主要是高强度、低挥发分的优质焦煤，这种煤蒙古和俄罗斯都无法替代。这意味着中国钢铁厂在配煤环节失去了关键的“骨架”原料，只能转向进口。完整报告中的图2和图3清晰地展示了这一冲击如何迅速拉平了澳大利亚煤和国内煤的价格。

更值得警惕的是结构性成本的变化。报告中提到，部分山西矿井的开采深度已达1-1.5公里，本身就处于成本曲线的高端。监管收紧后，安全投入的增加将进一步推高国内焦煤的边际成本。这意味著即使产能恢复，国内煤相对于进口煤的竞争力也在下降。

*对行业格局的含义：** 中国从“偶尔依赖进口”正在转向“结构性依赖进口”。过去四年，澳大利亚煤在中国进口中的占比一度降至7\%左右，但这次事件正在逆转这一趋势。

报告中最引人注目的数据点是：澳大利亚优质焦煤目前在中国港口的到岸价格，已经比国内煤便宜5美元/吨。这是过去四年多来首次出现的情况。

GS团队测算，自山西事故以来，中国CFR价格已上涨约40美元/吨。随着澳大利亚煤的竞争力显现，越来越多的现货正在转向中国市场。交易商反馈，一级和二级澳大利亚焦煤（包括PCI煤）正加速流向中国。

KC评论： 这意味着什么？过去几年，澳大利亚煤在中国市场几乎被边缘化，定价权主要掌握在蒙古和俄罗斯手中。但现在，澳大利亚正在重新成为中国的边际供应来源。当中国需要补充高强度焦煤时，澳大利亚是唯一能够大规模提供替代供应的来源。完整报告中的图1展示了这一价格套利窗口的打开过程，值得仔细研究。

GS的分析师进一步指出，随着澳大利亚煤重新进入中国市场，整体价格实现将得到改善。过去，部分矿商为了避免压低澳大利亚基准价格，选择直接向中国销售现货。但现在，随着中国和澳大利亚价格趋同，这一顾虑正在消除。

*对资产定价的含义：** 澳大利亚焦煤生产商的定价权正在修复。如果这一趋势持续，市场对澳大利亚矿业公司的估值模型可能需要上调——尤其是那些拥有优质低挥发分焦煤储量的公司。

印度是焦煤需求增长最快的市场。报告显示，印度目前每年进口约8000万吨焦煤，预计这一数字将增长到超过1.2亿吨。背后是印度钢铁需求的独特结构：基础设施和建筑占钢铁需求的64\%，远高于全球平均的52\%。

然而，GS团队在会议中捕捉到一个关键矛盾：印度虽然需求量大，但在价格形成中并不占据主导地位。

原因有三。第一，印度买家高度依赖长期合同，现货参与比例通常只有10\%左右。第二

[... middle omitted ...]

化——通过增加二级煤、半软煤和国内煤的比例，将优质焦煤的依赖度从90\%以上降至45-65\%。这一趋势正在结构性改变全球焦煤的需求结构，压低了对最高等级焦煤的边际需求。

报告揭示了一个容易被忽视的事实：澳大利亚的焦煤出口能力正在被多重因素侵蚀。

过去十年，澳大利亚焦煤出口已从约2.2亿吨下降至1.95亿吨。GS预计，到2030年还将再减少1000万吨。原因包括：矿山枯竭（如Oaky Creek的关闭）、剥离比上升、矿坑加深、以及储备质量下降。

成本方面，昆士兰州的运营成本已比2022年初高出30-40\%。柴油成本占运营支出的10-20\%，仅燃料通胀就增加了每吨8-12美元的成本。再加上昆士兰州的特许权使用费上调，多位与会者表达了对该地区竞争力下降的明显不满。

*对投资决策的含义：** 澳大利亚的供应弹性正在消失。新的矿山审批周期超过5年，现有的矿山延期也面临漫长审批。这意味着即使价格上涨，供应端也无法迅速响应。全球焦煤市场的供需平衡将更加脆弱，价格波动可能加剧。

\subsection{[R006] 摩根斯坦利：MS：中国5月成品油需求暴跌，电动车不是“冲击”，是“替代完成”}
{\small\color{refgray}归类：未被主题摘要引用}

摩根斯坦利：MS：中国5月成品油需求暴跌，电动车不是“冲击”，是“替代完成”

中国5月的成品油消费数据，发出了一个比市场预期更清晰的信号。

MS最新发布的研报显示，中国5月隐含汽油和柴油需求同比分别下降12\%和21\%，总成品油消费同比下降13.5\%。这个数字不是简单的“需求疲软”，而是结构性拐点的确认。报告明确指出，需求破坏的主要来源是公路运输和基建活动，而非航空或化工。

换句话说，中国石油需求的长期故事，正在从“增长放缓”切换为“存量替代”。

为什么这个判断比GDP增速或PMI数据更值得关注？因为成品油消费是经济活动最真实的“物理映射”——它不受库存周期、价格幻觉或统计口径的干扰。当汽油和柴油同时出现两位数下滑，且背后是电动车对燃油车的系统性替代时，投资者需要重新评估整个能源化工产业链的定价逻辑。

1. 柴油需求暴跌21\%，基建和物流的“真实温度”比想象更低

柴油是经济活动的“毛细血管”——它关联着卡车运输、工程机械、矿山开采和农业机械。5月21\%的同比降幅，是疫情以来最深的单月跌幅之一。MS将其归因于“基建和物流需求放缓”，但这背后可能还有更深层的结构变化。

传统上，柴油需求与固定资产投资高度相关。但这一次，数据似乎在说：即便基建投资名义增速没有大幅下滑，单位投资额的柴油消耗量正在下降。原因有二：一是新能源重卡和工程机械的渗透率在快速提升；二是物流行业的“集约化”趋势——更高效的路线规划和载具利用率降低了对柴油车的依赖。

KC评论： 柴油21\%的降幅不是“经济崩了”，而是“经济结构变了”。基建投资可能还在，但拉动的柴油消费正在被电动化和效率提升抵消。完整报告里对柴油库存和炼厂开工率的数据拆解，能帮你判断这是短期波动还是长期趋势。

2. 汽油需求下滑12\%，电动车不再是“边际冲击”，而是“主力替代”

汽油需求的下滑更值得细读。12\%的同比降幅，在历史上只有2020年初疫情封锁期间出现过。但这一次，没有封锁，没有大规模居家令。MS的判断是：高油价加速了消费者向电动车和公共交通的转移。

第一层是“替代速度”。2023年中国新能源乘用车渗透率突破40\%，2024年继续攀升。这意味着每新增100辆乘用车中，超过40辆是电动车。存量车中的燃油车替换速度也在加快——二手车市场燃油车价格持续承压，进一步降低了燃油车的持有意愿。

第二层是“行为锁定”。一旦消费者转向电动车，其加油行为几乎永久消失。这不是“今年少加一次油”的问题，而是“未来十年都不再加汽油”。因此，汽油需求的下滑不是线性的，而是带有“跳跃式”特征的——当渗透率突破某个临界点，替代速度会突然加速。

KC评论： 汽油需求12\%的降幅，本质上是“存量替代”而非“增量放缓”。报告中关于“高油价是否鼓励了消费者转向”的讨论，只是其中一个变量。更完整的分析需要看电动车保有量增速与汽油消费弹性系数的历史关系——这些在完整报告中有详细图表。

在汽油和柴油全线溃败的背景下，航空燃油5月同比增长5\%，石脑油消费基本持平。这似乎给成品油市场留下了一块“安全垫”。但MS的数据暗示，这个安全垫可能并不牢固。

航空燃油的韧性来自国内航空出行恢复和国际航线逐步放开。但问题是，航空燃油在中国成品油消费中的占比只有10\%左右，无法对冲公路运输的下滑。更重要的是，航空领域的电动化虽然尚未大规模展开，但氢能和可持续航空燃料（SAF）的试点正在推进——长期看，航空燃油也不是“免疫区域”。

石脑油的需求持平，反映的是化工产业链的稳定。但石脑油的下游——乙烯、丙烯等基础化工品——正面临产能过剩的压力。如果化工品价格持续低迷，石脑油的需求也难以独善其身。

MS在报告中特别提到一个矛盾现象：尽管炼厂大幅削减开工率，成品油库存仍然处于高位。这说明需求下降的速度超过了供应调整的

[... middle omitted ...]

库存是如何一步步累积的。这个趋势如果持续，可能触发更大规模的产能出清。

MS这份报告的核心贡献，是指出了“替代正在加速”这个方向。但它没有完全回答一个关键问题：这个替代速度何时会“二阶加速”——即从现在的两位数降幅，进一步扩大到更高的水平。

第一，电动车的成本下降速度是否快于燃油车的效率提升速度。如果电池成本继续下降，电动车的全生命周期成本优势将更加明显，替代速度会进一步加快。

第二，充电基础设施的瓶颈何时真正突破。目前充电桩的分布不均和充电时间仍然是制约因素。如果超充网络在 年大规模铺开，燃油车的最后一块“便利性优势”也将消失。

第三，政策是否会加速替代。碳达峰目标的推进、燃油车禁售时间表的明确化、以及可能的燃油税上调，都可能成为加速替代的“催化剂”。

这些变量，决定了成品油需求的下滑是“每年10\%”还是“每年20\%”。而这，直接关系到炼厂、油田服务和相关化工品的估值逻辑。

基于以上分析，投资者需要重新构建对中国能源化工产业链的观察框架。核心是三个维度：

\subsection{[R007] BARC：标普500目标价7800，但真正支撑上涨的不是估值，是盈利}
{\small\color{refgray}归类：风险偏好：美股上涨逻辑从估值驱动切换至盈利驱动，AI牛市未能推升美元}

BARC：标普500目标价7800，但真正支撑上涨的不是估值，是盈利

当市场还在为中东停火、油价回落、AI资本开支是否见顶而反复博弈时，BARC策略团队做出了一个看似矛盾的决定：上调标普500目标价，却同时下调了估值假设。

2026年底标普500目标价从7650上调至7800，2027年底目标价8800——这些数字本身并不惊人。真正值得关注的信号是，BARC将估值倍数从此前的约24倍下调至23.1倍，同时将2026年每股收益（EPS）预期从321美元大幅上调至337美元。

*这意味着什么？** BARC认为，推动美股继续上涨的发动机正在切换：从流动性驱动、估值扩张，转向盈利驱动、基本面支撑。如果这个判断成立，那么投资者需要重新审视自己的持仓结构——哪些公司真正有盈利可见性，哪些只是搭了估值泡沫的顺风车。

这份报告写于2026年6月19日，正值市场经历了一轮从地缘政治冲击到快速修复的完整周期。BARC的结论是：美股的牛市基础仍然完整，但支撑逻辑已经变了。

1. 这轮上涨的底层逻辑，正在从“估值讲故事”转向“盈利证明自己”

2026年上半年的美股走势，像一部浓缩的风险教育片。年初市场已经定价了大量好消息：全球宏观企稳、AI资本开支成为不可否认的顺风、市场领导力从超大型科技股向周期股和小盘股扩散。但高估值意味着容错空间极窄。

2月，软件行业被AI颠覆的担忧引发抛售，跌幅抹去了该板块自2022年以来相对于标普500的全部超额收益。随后，对高杠杆软件公司的担忧蔓延至私募信贷市场，进而引发对长久期成长资产的全面重估。紧接着，中东战争爆发，石油价格飙升，市场开始担忧滞胀。

BARC的策略是“看穿噪音”。当停火协议传出后，投资者迅速回归那些盈利可见性最强的板块——大科技、成长股和动量股重新领涨。2026年一季度财报季表现强劲，科技板块EPS增速创下2010年以来最佳水平，全年盈利预期上修速度创历史纪录。

*但关键转折点在于：** BARC认为，随着美联储降息路径收窄、市场仓位已不再能轻易吸收负面冲击，盈利和AI资本开支的可见性必须承担更多“支撑市场”的任务。换言之，过去两年靠“美联储会救市”的信仰推高的估值，现在需要真实盈利来兑现。

KC评论： BARC的潜台词很明确——别再指望美联储给你送钱了。利率环境正在从“宽松预期”转向“高利率常态化”，这意味着估值扩张的红利基本耗尽。接下来市场会非常挑剔：只有那些能用盈利增长证明自己配得上当前估值的公司，才能继续获得资金青睐。完整报告里有一张关键的“收益率与股权相关性”图表，显示两者关系已转为显著负相关，这意味着利率上行对估值的压制正在加速——这张图值得仔细看。

2. 盈利上修的三个支柱：科技、再通胀、工业，但消费是最大拖累

BARC将2026年标普500 EPS预期从321美元上调至337美元，略低于市场共识的341美元，隐含同比增长约21\%。这个上修幅度不小，但更值得关注的是上修的结构——它揭示了盈利增长正在从哪些行业来，又将从哪些行业走。

*第一个支柱：科技。** 2026年一季度财报季的强劲表现，证实了大科技和TMT板块仍在将AI相关需求转化为盈利超预期。全年盈利上修速度远超10年均值。半导体和设备板块的2026年EPS增长预期从年初的66\%飙升至127\%，科技硬件从14\%升至39\%。软件板块则从18\%降至13\%，反映了AI颠覆的负面影响。

*第二个支柱：再通胀。** 2026年全年温和的通胀压力将支撑名义收入增长，尤其有利于那些有经营杠杆或商品敞口的行业。

*第三个支柱：工业。** 工业活动数据持续向好：工业生产上升、ISM制造业PMI终于回到扩张区间、耐用品订单和产出均显示强劲动能。BARC经济学家预测，2027年工业生产背景将进一步改善。



[... middle omitted ...]

首次给出了2027年标普500 EPS预期：389美元，同比增长约15\%，低于市场共识的398美元。这个数字隐含的判断是：2027年盈利增速将自然放缓，但不会出现断崖式下跌。

框架假设：消费在同比基础上部分复苏，工业活动继续走强（因为之前紧缩政策的滞后效应消退），但名义定价能力正常化（尤其是当去通胀力量扩大时）。乐观和悲观情景分别为397美元和380美元。

*这意味着什么？** BARC认为，2027年市场将面临一个“盈利增速放缓但仍有正增长”的环境。对于股市而言，这通常意味着估值承压，但盈利仍能提供底部支撑。关键在于，哪些行业能在增速放缓的环境中保持相对优势。

4. 板块轮动信号：金融和医疗从“看好”降级为“中性”，消费仍被看空

*金融板块从“看好”降级为“中性”。** 理由很直接：之前看多金融的逻辑——银行驱动的盈利上修——没有兑现。私募信贷担忧、支付领域的监管风险、以及AI颠覆的威胁，共同抵消了银行的正面盈利修正。BARC坦承“看多金融的论点没有走通”。

\subsection{[R008] DB：机构外汇管理正从“被动不管理”转向“主动总组合”}
{\small\color{refgray}归类：外汇管理与机构配置：从被动不管理转向主动总组合}

当一家大型养老金基金的管理者被问及如何管理汇率风险时，最常见的回答是“我们没有专门做这件事，因为它分散了我们对核心资产的关注”。这种回答正在被重新审视。

DB近期发布了一份关于外汇风险管理从理论到实践的深度报告。报告的核心判断并非某个汇率走势的预测，而是一个更根本的转变：机构投资者正在将汇率风险管理从“后台被动合规”提升为“总组合层面的主动策略”。这家投行在与超过35家机构客户讨论后得出结论——不管理汇率风险本身就是一种风险决策，且往往是最低效的那种。

为什么是现在？三个因素正在汇聚：第一，荷兰养老金改革等区域监管变化正倒逼机构重新审视风险敞口；第二，全球资产配置中非本币资产的比重持续上升，汇率波动对总回报的影响已无法忽略；第三，机构对“总组合方法”（TPA）的兴趣从理论探讨转向了如何落地执行。

这份报告的价值不在于给出具体的做多或做空建议，而在于提供了一套机构可以照搬的思考框架。它回答了“管不管”、“管多少”、“用什么工具管”以及“管了之后如何应对流动性压力”这四个递进的问题。

报告开篇就直指一个长期存在的行业惯例：汇率风险往往被忽视，或者被分割到各个资产类别中单独管理。债券组合可能做了对冲，股票组合可能完全暴露，而私募市场的汇率风险则基本没有被系统性地处理。

这种“碎片化管理”的后果是什么？DB指出，选择不对冲本身就是一种风险决策。如果这个决定是被动做出的——即因为“没人负责”或者“太复杂”——那么它对组合结果的影响往往是不对称的。当汇率波动恰好与核心资产回报负相关时，组合会承受双重打击。

从区域分布来看，亚太地区的机构在讨论中对汇率管理的关注度最高，而北美地区相对最低。这背后的逻辑不难理解：亚太机构持有大量美元、欧元资产，而本币（如日元、韩元、新加坡元）的波动对这些机构的实际购买力影响巨大。北美机构则因为美元的国际地位，往往低估了非美元敞口的风险。

KC评论： 这里的核心洞察是“不决策本身就是决策”。很多投资决策者把“不对冲”当作一个中性选项，但报告提醒，在总组合框架下，不对冲意味着把汇率风险的定价权完全交给了市场。读者如果管理的是跨境资产组合，第一个需要回答的问题是：我的汇率敞口是主动选择的，还是因为没人管而被动存在的？

这是报告中最具洞察力的论点之一。在机构客户的讨论中，一个根本性的争议始终存在：管理汇率风险到底是为了增厚收益，还是为了控制波动？

DB的立场非常明确——汇率管理应该服务于组合的稳定性，而不是追求汇兑收益。报告指出，核心资产的任务是产生回报，而汇率管理的作用是“在保持组合回报的同时管理风险”。换句话说，不要指望靠做多日元或做空欧元来赚钱，那是外汇交易员的活。机构投资者的汇率管理应该是一个“保险”功能，而不是一个“收益”来源。

这个结论对很多机构来说并不容易接受。因为在实际操作中，当某一货币大幅升值时，完全对冲的机构会错失这部分收益，而部分对冲的机构则可能被质疑“为什么没有多赚一点”。但报告用了一个简洁的框架来化解这个矛盾：把收益来源归到核心资产，把汇率管理定位为运营效率工具。

KC评论： 这个框架对投资委员会有直接的参考价值。它意味着，评价汇率管理团队的标准不应该看“他们赚了多少汇兑收益”，而应该看“他们在多大程度上降低了组合的波动率”。如果一家机构的绩效评估体系把汇率损益和股票回报混在一起算，那么它很难真正执行一套纪律性的对冲策略。这份报告对如何设计这样的评估体系有更详细的拆解，值得完整阅读。

机构投资者不是自由交易的。大多数基金都有明确的业绩基准，比如MSCI全球指数，而这个指数通常是不对冲的。这意味着，如果一家基金大幅对冲了汇率风险，而基准没有对冲，那么在汇率大幅波动时，基金的表现会显著偏离基准。

这种“跟踪误差”约束是很多机构不敢大胆管理汇率风险的根本原

[... middle omitted ...]

（即利率差）中。当对冲期限较长或利率差较大时，这个成本可能相当可观。期权则提供了非线性回报，机构可以通过卖出期权来收取权利金，从而部分覆盖对冲成本。

报告特别提到，那些必须维持完全对冲的基金（如某些荷兰养老金），正在利用期权来引入灵活性。例如，通过买入价外期权来保护极端风险，同时用卖出期权来补贴保费。这种策略的核心在于“在风险管理和成本之间找到平衡点”。

这是一个容易被忽视但极其重要的维度。报告指出，机构需要从“直接持仓”转向“经济敞口”的视角来看待汇率风险。

举例来说，如果一家基金持有一家美国上市公司的股票，而这家公司的大部分收入来自日本，那么这只基金实际上承担了日元敞口，尽管它持有的资产是以美元计价的。这种间接敞口无法通过简单的对冲比例计算来捕捉。

报告坦承，目前没有简单的方法来精确衡量这种经济敞口。一个可行的思路是：逐个分析公司年报和季报，估算其收入、成本和利润的货币分布。这听起来繁琐，但对于大型基金来说，前20大持仓往往就决定了绝大部分的间接汇率风险。

\subsection{[R009] NOM：电池行业已进入“三国分治”阶段，机器人电池是下一个价值高地}
{\small\color{refgray}归类：新能源车与电池：订单分化加速，电池行业进入三国分治阶段}

NOM：电池行业已进入“三国分治”阶段，机器人电池是下一个价值高地

全球电池产业正在经历一场静悄悄但极其深刻的分化。当多数市场参与者还在用“产能过剩”和“价格战”来概括行业现状时，NOM最新发布的全球电池深度报告给出了一个截然不同的判断：行业已经不再是单一维度的竞争，中国、韩国、日本三大电池阵营正在走向完全不同的战略路径，而机器人电池——这个目前仅占极小份额的细分市场——可能成为未来几年最被低估的价值增长点。

这份报告的核心洞察并不复杂，但其含义却值得每一位关注先进制造和能源转型的决策者深思：全球电池需求仍在增长，但增长的结构正在发生质变。储能系统（ESS）电池将以17\%的年复合增速成为增长最快的板块，到2030年达到926GWh；而电动汽车电池的增速将放缓至11\%。更重要的是，中国企业凭借LFP电池和一体化供应链，已经拿下了全球78\%的电动汽车电池和80\%的储能电池市场份额。韩国和日本企业不再试图在规模上与中国正面竞争，而是分别转向了“高价值+本地化”和“AI数据中心”两个截然不同的赛道。

NOM报告中最引人注目的数字，是中国电池企业在全球市场的份额。2026年，中国在全球电动汽车电池和储能电池市场的份额预计分别达到78\%和80\%。这个数字之所以惊人，不仅在于其绝对高度，更在于它是在美国和欧洲持续加征关税、设置本地化壁垒的背景下实现的。

中国电池的统治力根植于LFP（磷酸铁锂）化学体系的全面胜利。目前LFP已经占全球电池市场的61\%，而中国几乎完全掌控了这条供应链。从锂资源加工到正极材料、隔膜、电解液，中国企业在每一个环节都拥有规模优势和成本优势。NOM报告指出，即使在海外市场，中国电池厂商凭借成本领先和供应链整合能力，依然占据了58\%的份额。

KC评论： 这意味着，对于全球车企和储能项目开发商而言，试图完全“去中国化”的电池供应链策略，在短期内不仅成本高昂，而且可能面临技术和交付能力的巨大缺口。中国电池的统治力不是政策补贴的产物，而是十年持续投入形成的系统级优势。完整报告中有详细的成本对比图表和各国LFP产能地图，这些数据会进一步说明为什么“替代中国电池”是一个比多数人想象中更困难的工程问题。

值得注意的是，NOM认为中国在下一代电池技术上依然保持领先。钠离子电池预计到2030年将占据储能电池需求的15-20\%，而中国同样掌控着钠离子电池的供应链。半固态电池可能在未来两年内率先在中国实现商业化。这些技术进展将进一步巩固中国电池企业的护城河。

韩国电池企业的战略转型，是这份报告中最值得关注的竞争格局变化之一。面对中国在规模上的绝对优势，韩国企业不再试图在成本上正面竞争，而是系统性地向两个高壁垒领域转移：高端电动汽车电池和美国储能市场。

在高端电动汽车领域，韩国企业保留了竞争优势。高镍化学体系（NCM/NCA）在长续航车型中仍有不可替代性，特斯拉青睐的4680圆柱电池也是韩国企业的强项。硅负极技术同样是韩国企业重点布局的方向。但真正形成壁垒的是美国和欧洲的本地化要求——关税壁垒和本地化生产要求正在为韩国企业创造出一个“有保护溢价”的市场空间。

NOM报告特别指出，韩国电池企业在美国储能市场的渗透正在加速。美国AI数据中心电力需求、可再生能源并网和电网现代化正在推动储能需求快速增长，预计到2030年美国储能电池需求将达到194GWh。在这个市场中，中国电池受到关税和地缘政治因素的制约，韩国企业因此获得了显著的定价权和增长空间。

KC评论： 韩国电池企业的故事，本质上是一个“从追求规模转向追求定价权”的战略案例。对于投资者而言，关键问题在于：美国储能市场的本地化溢价能持续多久？如果中美贸易关系缓和，韩国企业的护城河是否会被削弱？完整报告中对韩国三大电池企业的产能规划和美国项目储备有详细拆解，这些

[... middle omitted ...]

增长极快。NOM估计，松下在全球AIDC BBU市场中占据了60-70\%的份额，并且设定了到2029财年实现1万亿日元收入的目标，同时要求整个电池业务的ROIC达到20\%。这个数字对于习惯了电池行业“薄利多销”思维的观察者来说，是一个令人意外的信号。

日本企业的选择背后有一个核心逻辑：在AI数据中心这个场景中，电池的性能要求远高于电动汽车和普通储能。高功率密度、高倍率放电能力、轻量化设计和先进的热管理，这些技术参数恰恰是日本企业在材料科学和精密制造领域的传统优势。在这个细分市场中，价格不是首要竞争维度，可靠性和性能才是。

4. 机器人电池：被忽视的高价值蓝海，单价是电动汽车电池的三倍

NOM报告中最具前瞻性的判断，可能来自对机器人电池市场的分析。人形机器人的商业化正在加速，而机器人对电池的需求与电动汽车完全不同。机器人电池需要更高的功率密度、更高的C-rate（充放电倍率）能力、更轻的重量和更先进的热管理系统。这些要求使得机器人电池的单价达到电动汽车电池的三倍。

\subsection{[R010] GS：炼油利润的韧性比原油价格更值得关注}
{\small\color{refgray}归类：能源与大宗：炼油利润韧性优于原油，铜矿股定价过度悲观}

原油价格在美伊临时和平协议宣布后暴跌超过10美元，但炼油产品的利润率降幅却小得多。GS全球炼油产品利润率指数目前仍处于战前水平的两倍。

这是一个被市场低估的信号。当所有人都在盯着布伦特原油的波动时，真正影响产业链利润结构的关键变量正在发生变化。GS最新发布的原油研报给出了一个核心判断：炼油产品利润率的下降风险远小于原油价格，且结构性支撑因素将使其在2027年仍维持高位。

这份报告的价值不在于预测油价的方向——那是几乎不可能精准完成的任务——而在于揭示了能源产业链中利润分配的结构性变化。对于关注能源资产定价、炼化板块投资以及通胀传导的决策者而言，这可能是比原油价格本身更重要的观察维度。

1. 霍尔木兹海峡的重新开放对炼油利润的冲击远小于对原油的冲击

市场普遍认为，一旦霍尔木兹海峡通行恢复正常，所有与地缘溢价相关的资产价格都会同步回落。但GS的数据显示，实际情况并非如此。

自美伊临时和平协议宣布以来，原油价格下跌超过10美元，而炼油产品利润率的降幅明显更小。GS的全球炼油产品利润率指数仍维持在战前水平的两倍。这一差异背后是两种完全不同的供需逻辑。

原油价格受制于全球供需平衡和地缘政治溢价的快速消退，而炼油产品利润率则受到更具体、更刚性的约束——炼厂产能利用率、产品库存水平以及不同馏分油的生产调配。这些因素不会因为一份和平协议就迅速改变。

KC评论： 这组数据指向一个容易被忽视的事实：即便地缘风险溢价消退，炼油环节的利润也不会同步回归“正常”。对于持有炼化资产或关注相关板块的投资者而言，这意味着需要重新评估估值框架中的利润中枢假设。完整报告中附有详细的全球炼油产品利润率指数走势图，展示了与原油价格的背离程度，值得仔细研究。

GS看多炼油利润率的核心理由是“更长时间的低库存”。无论是汽油还是柴油，美国的库存水平都低于 年的季节性区间。更关键的是，GS预计柴油和汽油库存在未来1-4个月内将继续下降，直到2027年下半年才会逐步恢复到2025年的季节性水平。

汽油库存尤其紧张。过去四个月，炼厂优先生产柴油和航空燃油，以牺牲汽油产量为代价。这种生产结构的调整不是短期行为，而是对利润信号的理性反应。

与此同时，尽管停火协议宣布后波斯湾炼厂的非计划停产规模减少了一半，但仍有约130万桶/天的产能处于非计划维修状态。GS指出，这些产能可能需要数月时间才能修复。

KC评论： 这里的核心变量是“时间”。130万桶/天的停产产能不是开关一按就能恢复的，维修周期以月计。这意味着即便地缘政治风险消退，供应侧的恢复速度也会显著慢于市场预期。完整报告中的Exhibit 6展示了波斯湾炼厂停产的具体恢复路径，对于判断未来几个季度的供应节奏至关重要。

3. 2027年的结构性支撑来自中印炼能扩张延迟和俄罗斯炼厂持续受袭

GS对2027年炼油利润率的乐观预期并非基于短期供需扰动，而是建立在三个结构性因素之上。

第一，中国和印度的炼油产能扩张出现明显延迟。这两个全球最大的炼油产能增长来源正在放缓扩张步伐，这意味着全球炼油产能的增长速度低于此前预期。

第二，俄罗斯炼厂持续遭受袭击。即便不考虑地缘政治谈判的进展，俄罗斯炼油产能的恢复也面临技术和安全层面的实际困难。

第三，全球炼油产能利用率将维持在接近历史高位的水平。GS预计，即便到2027年，全球炼油产能利用率仍将处于历史高位区间。

这三个因素叠加，意味着炼油利润率的支撑不是短期现象，而是一个中期结构性特征。GS预计2027年美国柴油利润率为38美元/桶，欧洲为25美元/桶，虽然较2026年有所回落，但仍显著高于2025年的水平。

这是整份报告中最具投资含义的判断。GS通过情景分析量化了这种不对称性。

在悲观情景下——假设波斯湾出口在7月初恢复正常、原油产量超预期、需

[... middle omitted ...]

保护更强，上行弹性在某些品种上甚至更大。

KC评论： 这种不对称性是这份报告最值得关注的投资含义。对于持有原油多头头寸的投资者，下行风险显著大于炼油利润；而对于关注炼化板块的投资者，利润中枢的下行空间相对有限。完整报告中的Exhibit 8和Exhibit 9详细展示了不同情景下的利润路径，是理解这种不对称性的关键图表。

GS的预测框架中有一个关键假设：炼油产品需求将在未来一个季度内大幅反弹，2027年仅有温和的需求“疤痕”。这个假设决定了利润率预测的可靠性。

如果全球经济放缓超预期，或者能源转型加速导致结构性需求下降，炼油利润率的支撑逻辑将受到挑战。GS自己也承认，需求损失可能比基准假设更“粘性”——这也是其悲观情景的核心变量之一。

另一个未被充分讨论的问题是：炼油产能利用率维持高位是否意味着新产能投资的加速？如果炼油利润率持续高于历史均值，资本开支周期可能提前启动，这将在更长期内压制利润率。GS的报告提及了产能扩张延迟，但没有详细讨论产能投资周期的转折点。

\subsection{[R011] NOM：中国AI投资的真正引擎不是BAT，是“国家资本矩阵”}
{\small\color{refgray}归类：AI/算力：需求三年翻五倍，MLCC与CoWoS成为新瓶颈}

NOM：中国AI投资的真正引擎不是BAT，是“国家资本矩阵”

理解中国AI产业的融资结构，比理解其技术路线图更为关键。NOM在最新发布的亚洲宏观研究报告中，系统拆解了中国AI资本开支的资金来源。其核心判断是：中国AI投资的驱动力并非市场自发行为，而是一个由国家资本主导、多层级联动的“资金矩阵”。这与美国的私人资本主导模式形成了根本性差异。

这份报告的价值不在于罗列数字，而在于揭示了资金流向背后的政策逻辑：中国正在用财政工具和国有资本，来对冲外部技术封锁和内部投资下行压力。对于关注中国科技产业、资产配置以及地缘风险的读者而言，理解这个资金矩阵的运作方式，是判断AI赛道未来走向的基本前提。

报告指出，中国AI资本开支的规模与美国仍有量级差距，但其资金来源的“结构性”差异，可能比资金总量本身更重要。美国AI资本开支的绝大部分由超大规模云服务商驱动，而中国则通过中央与地方的多层资金工具，将“国家战略”转化为“实际投资”。

国家集成电路产业投资基金，俗称“大基金”，是中国半导体和AI领域最核心的国家级投资工具。NOM报告显示，大基金已运行三期，注册资本从第一期的1390亿元，跃升至第三期的3440亿元，累计规模接近7000亿元。这三期投资的演进路径，清晰地反映了中国产业政策的阶段性重心：第一期聚焦芯片制造、封测等下游环节；第二期向上游设备、材料等高壁垒领域延伸；第三期则全面覆盖AI产业链，包括存储芯片、光刻设备，甚至参与了AI初创公司DeepSeek的早期融资谈判。

KC评论： 大基金的投资路径不是随机的，它反映了一个清晰的“补链”逻辑。从下游制造到上游设备，再到AI全链条，每一步都在试图填补美国制裁留下的空白。对于投资者而言，大基金三期投资的领域，大概率就是未来几年政策扶持强度最高的“安全区”，但也意味着竞争格局将更加拥挤。

大基金的投资规模固然庞大，但它只是国家资本矩阵中的一个节点。当它与其他财政工具叠加时，其杠杆效应和信号意义才真正显现。

报告揭示了一个关键的政策转向：超长期特别国债的资金投向正在发生结构性变化。2025年，8000亿元“两重”建设资金主要流向交通、能源等传统基础设施。但到了2026年，情况发生了显著变化。在4月公布的第二批2170亿元资金中，人工智能被明确列为重点支持领域。这标志着，曾经为高铁、高速公路提供资金的财政工具，正在正式转向AI相关产业。

NOM还援引彭博社6月的一则报道称，中国政府正在起草一项为期五年、总额高达2万亿元的计划，旨在建立一个全国性的数据中心网络。该网络将依赖国产供应商提供至少80\%的核心技术，包括华为的AI芯片，并由中国移动、中国电信等国有电信巨头运营。若计入配套电网投资，总项目支出可能达到至少5万亿元。资金将主要来自超长期特别国债和战略性产业国有基金。

KC评论： 2万亿数据中心计划如果属实，其意义不亚于当年的“村村通”工程。它试图构建一个从芯片到算力、再到运营的完全国产化闭环。但这也提出了一个关键问题：如此庞大的算力供给，能否被市场需求有效消化？报告没有详细展开需求侧的测算，这恰恰是完整报告中最值得深读的部分。

除了直接的财政拨款，中国还动用了“准财政”工具。报告提到的“新政策融资工具”于2025年9月启动，初始额度5000亿元，2026年又追加了8000亿元。与 年和2022年的前两轮类似工具不同，本轮资金明确优先投向数字经济、AI和消费基础设施。国家开发银行和进出口银行分别将37.5\%和40\%的资金配置给了数字经济和AI领域。

更关键的是，这些资金作为项目资本金，能产生巨大的投资杠杆效应。据国家发改委估算，初始的5000亿元资金预计将撬动超过7万亿元的总项目投资。这意味着，国家投入的每一元钱，理论上可以带动14元的社会总投资。

与中央政府的

[... middle omitted ...]

础的地方政府仍能集中资源进行“重点突击”。

KC评论： 地方政府的AI投资呈现“冰火两重天”。大部分地区受困于财政压力，无力大规模投入；少数“幸运儿”则能抓住全球产业周期，获得中央和地方的双重加持。这提示我们，未来中国AI产业的区域分化可能会进一步加剧，投资机会将高度集中在少数几个“政策+产业”共振的城市。

在间接融资之外，股权市场也在发挥独特的“定向输血”作用。报告指出，尽管A股整体IPO自2024年初以来基本暂停，但科技公司仍被允许在科创板上市融资。过去一年，国产GPU设计公司摩尔线程（80亿元）、壁仞科技（39亿元）以及先进封装龙头盛合晶微（48亿元）均已完成IPO。

与此同时，香港市场也通过放松IPO规则来吸引内地科技公司。PCB制造商胜宏科技以230亿港元融资成为今年迄今香港最大IPO；AI芯片设计公司澜起科技完成了80亿港元的双重上市；智谱AI和MiniMax则分别融资44亿港元和55亿港元。科创板与港股市场，共同构成了中国AI产业的“资本双通道”。

\subsection{[R012] HSBC：中国AI算力需求三年翻五倍，这两家公司最受益}
{\small\color{refgray}归类：AI/算力：需求三年翻五倍，MLCC与CoWoS成为新瓶颈}

当大多数投资者还在为AI应用能否落地而争论时，算力基础设施的供需缺口已经给出了最直接的答案。

HSBC证券最新发布的研报给出了一个清晰且大胆的判断：中国AI算力需求将在2028年达到2025年的5倍，而供给端受制于建设周期和GPU产能，将持续处于紧张状态。这意味着，算力租赁价格不仅不会下降，反而有望继续上行。

这份报告的核心逻辑并不复杂：token消耗量在爆发式增长，中国超大规模云服务商（CSP）的资本开支在2026年将突破7000亿元人民币，接近翻倍。而新建一个AI数据中心通常需要18-24个月，供给缺口短期内无法弥合。

在这轮算力军备竞赛中，HSBC明确看好两类公司：一类是拥有大规模AIDC产能的运营商，另一类是数据中心交换机供应商。报告首次覆盖了润泽科技（Range）和锐捷网络（Ruijie），均给予买入评级。

KC评论： HSBC的核心判断与市场普遍预期的“算力过剩论”形成鲜明反差。当很多人担心DeepSeek等高效模型会降低算力需求时，报告用token消耗量的实际数据证明，模型效率提升反而刺激了更大的使用量。这是一个典型的“杰文斯悖论”——效率提升带来更多需求，而非减少。

HSBC估算，中国智能算力需求将从2025年的约0.7GW增长至2028年的约5.4GW，年复合增长率高达74\%。这个数字的支撑点，是token消耗量的真实数据。

根据国家数据局的数据，2026年3月中国日均token消耗量达到140万亿，较2025年12月增长40\%。IDC的预测更为激进，全球智能体token使用量将从2025年的0.0005 peta增长至2030年的152,667 peta。

这背后的驱动力是生成式AI渗透率的快速提升。字节跳动、阿里巴巴、腾讯三大巨头的合并资本开支将从2025年的3460亿元翻倍至2026年的超过7000亿元。即便如此，这个数字仍远低于美国同行的投入强度，意味着增长空间依然巨大。

HSBC据此推算，中国第三方数据中心市场将在2028年达到1650亿元，其中AIDC市场将以78\%的年复合增速增长至750亿元。

算力需求的爆发只是一面。真正让这个投资逻辑成立的是供给端的刚性约束。

HSBC报告指出，Nvidia H100的一年期租赁价格已从12元/GPU/小时上涨至17元/GPU/小时，涨幅约40\%。这不是个别现象。2026年以来，腾讯云、阿里云、百度云、Ucloud等相继宣布AI算力产品涨价5\%-34\%，Google Cloud和Amazon AWS也在北美和欧洲市场同步提价。

供给端出现瓶颈的原因有三个：GPU产能受限、AIDC建设周期长、以及SuperPOD架构带来的技术壁垒提升。

一个典型的AIDC项目从建设到上架需要18-24个月。即便现在所有玩家都在加速扩张，短期内能释放的有效产能依然有限。HSBC统计了主要第三方IDC供应商的在运容量和规划，发现头部企业的储备容量远无法满足 年的需求增量。

KC评论： 这里有一个容易被忽略的细节：GPU租赁（GPUaaS）和传统IDC租赁是两种完全不同的生意。GPUaaS是运营支出，企业可以按月付费，而传统IDC需要长期锁定。这意味着GPUaaS的定价弹性更大，涨价传导更快。完整报告中有一张商业模式对比表，清晰展示了五类数据中心业务在定价模式、服务内容和目标市场上的差异，值得细看。

HSBC报告提出了一个关键的技术变量：SuperPOD集群架构。

SuperPOD本质上是将多个GPU通过高带宽互联整合到一个机柜内，实现大规模并行训练。这种架构对数据中心提出了全新的要求：单机柜功耗从传统CPU机柜的4-12kW飙升至120kW以上，同时要求更高的楼层承重、层高、液冷能力和网络架构。

技术门槛的提升，直接带来了运营商价值主

[... middle omitted ...]

动力：

第一，更大的GPU集群需要更多网络层级。传统数据中心网络通常只有两到三层，但万卡甚至十万卡集群需要更复杂的网络拓扑，交换机数量随之增加。

第二，端口速率升级带来单机价格提升。从100G到400G再到800G，交换机端口速率每提升一代，单机价格往往翻倍。

第三，SuperPOD架构推动了Scale-up交换机的渗透。这类交换机用于连接机柜内的GPU，是全新的增量市场。随着SuperPOD渗透率上升，这一细分市场将成为交换机厂商的重要增长极。

HSBC预计，锐捷网络在中国数据中心交换机市场的份额将从2025年的21\%提升至2028年的24\%。

尽管HSBC的逻辑链条完整，但这份报告在两个问题上留下了悬念。

第一个问题是：国产替代对算力供给格局的影响有多大？报告主要基于Nvidia GPU的供需框架进行推演，但华为昇腾等国产芯片的产能爬坡速度、性能差距和生态成熟度，都可能改变供给曲线。如果国产GPU在 年实现规模替代，算力租赁价格的上涨空间可能会被压缩。

\subsection{[R013] 摩根斯坦利：MS：AI硬件下一个瓶颈不是GPU，是MLCC}
{\small\color{refgray}归类：AI/算力：需求三年翻五倍，MLCC与CoWoS成为新瓶颈}

当市场还在为英伟达GPU的供应周期和ASIC的替代叙事争论不休时，MS最新发布的全球科技网络研讨会报告指向了一个更隐蔽、但可能更早到来的瓶颈——MLCC（多层陶瓷电容）。

这不是一个关于“缺芯”的老故事。这是AI服务器功耗架构演进带来的被动元件质变。报告的核心判断是：MLCC正在从“跟着服务器出货量走”的被动元件，变成AI服务器性能升级的主动瓶颈。到2027年，AI服务器对MLCC的需求将逼近10亿美元，而全球前五大供应商的产能布局和产品升级节奏，将决定整个AI硬件供应链的交付上限。

这个判断的背后，是三个正在同时发生的结构性变化：AI服务器功耗架构要求MLCC从“外围滤波”走向“近CPU/GPU嵌入”，中国AI芯片在推理侧的成本竞争力正在倒逼全球供应链重新定价，以及MLCC行业本身已经高度寡头化带来的供给刚性。

KC评论： 这份报告最值得读的部分，不是MLCC需求预测本身，而是它如何把MLCC、中国AI芯片、先进封装和功率半导体放在同一个框架里分析。四者之间的联动关系，才是理解AI硬件供应链当前状态的关键。完整报告里有大量图表支撑这些联动判断，这里只提炼核心逻辑。

传统服务器中，MLCC是标准化的通用元件，需求增长主要跟着服务器出货量走。但AI服务器的功耗和电流密度已经让传统方案失效。MS分析师指出，新架构要求MLCC具备更高的电容值、更低的等效串联电感（ESL），并且要能嵌入到更靠近CPU和GPU的位置。

这意味着三件事：单台AI服务器的MLCC用量和价值量都在提升；对供应商的技术能力要求跳升了一个台阶；产能规划不再是“按历史增速线性外推”，而是需要提前两到三年锁定先进产能。

报告估算，AI服务器MLCC需求将从2025年的约4亿美元增长到2027年的约10亿美元。这个增速远快于传统服务器MLCC市场。但问题在于，全球MLCC市场87\%的份额集中在五家供应商手中——村田（Murata）、三星电机（SEMCO）、国巨、TDK和太阳诱电。前两家被MS评为超配（OW）。

当一个高度寡头的行业遇到一个突然加速的需求曲线，结果通常是：价格弹性下降，交付周期拉长，下游客户被迫接受更高的库存水位。

KC评论： 这里有一个容易被忽略的点。报告提到AI服务器MLCC需求增长，但没有详细展开的是：这10亿美元需求中，有多少是现有产能可以满足的，有多少需要新建产能。如果前五大供应商的扩产节奏跟不上，MLCC可能成为继CoWoS之后的下一个“交期拉长、价格跳升”的环节。完整报告里有各家供应商的产能规划对比，值得细看。

MS用了相当篇幅分析中国AI芯片生态。结论有些反直觉：在推理场景下，国产芯片的每token成本已经可以和英伟达的中国特供版处理器竞争。

报告将中国AI芯片厂商梳理为“十大金刚”，重点覆盖了寒武纪（超配）、沐曦（等权重）和壁仞科技（超配）。核心论据有两个：一是国产芯片在FP8/FP16精度下的推理性能/成本比优于英伟达的合规产品；二是中国在基础设施层面的优势——电力、网络、数据中心建设成本——进一步拉低了推理的TCO。

报告给出的数据很具体：寒武纪MLU590在DeepSeek R1模型下的推理性能/成本比，已经接近甚至超过英伟达同类产品。而字节跳动豆包/火山引擎的token量激增，也从需求侧验证了中国市场对AI推理的强劲需求。

但这里有一个关键伏笔：国产芯片的推理性能优势，很大程度上建立在较低的定价和基础设施成本上。如果英伟达调整中国市场的产品策略和定价，或者美国出口管制进一步收紧先进制程代工，这个优势的可持续性需要重新评估。

KC评论： MS对中国AI芯片的覆盖，最有价值的部分不是个股推荐，而是他们搭建的“九因素比较框架”——从芯片、系统到基础设施，全面对比中美AI算力差距。报告认为中国在

[... middle omitted ...]

远超市场预期。

AWS的Trainium同样在加速。Trainium3预计2025年开始小批量出货，2026年达到160万颗，随后Trainium4在 年接力。加上Meta的MTIA和特斯拉的Dojo，全球主要CSP的定制芯片计划已经不再是“备选方案”，而是明确的长期路线图。

这对英伟达意味着什么？报告没有直接下结论，但数据本身已经说明问题：当CSP的自研芯片出货量从百万级向千万级迈进时，英伟达在AI训练市场的份额天花板将不再由技术决定，而由客户的自研意愿和代工产能决定。

报告的另一条主线是功率半导体。MS观察到，全球主要功率半导体公司的资本开支已经连续两年下降，而需求端——尤其是工业自动化和中国新能源汽车——正在边际改善。

工业自动化企业的收入在2026年一季度同比增长21\%，这是一个信号。中国新能源汽车批发量虽然年初至今偏弱，但边际上已经转正。更重要的是，功率半导体的供给端没有跟上：产能扩张预计在 年放缓，而汽车和工业应用合计占离散半导体需求的70\%左右。

\subsection{[R014] GS：美国变压器缺的不是产能，是中国速度}
{\small\color{refgray}归类：供应链与运价：高运价从事件冲击转向供给约束，中国变压器交付速度成优势 / 稀土磁材与电力设备：中国之外产能规划漂亮但产出存疑，变压器交付速度成差异化优势}

这份来自GS的最新变压器出口追踪报告，在五月数据中揭示了一个被市场低估的信号：中国变压器总出口增速从4月的27\%同比飙升至5月的72\%。但真正值得关注的，不是这个数字本身，而是数字背后的结构性错配——美国本土变压器供需缺口高达72\%，而中国供应商的交付速度是本土企业的四倍。

这不是一个简单的“中国出口强劲”的故事。这是一场关于全球电力设备供应链重构中，速度与规模如何重新定义竞争格局的博弈。

GS在这份报告中给出了明确的判断：美国电力变压器需求与本土供应的缺口，将从当前的72\%收窄至2028年的57\%。这个收窄幅度意味着什么？意味着即便考虑了所有已公布的美国本土扩产计划，短缺仍将持续存在。非传统供应商，包括中国企业，仍有明确的进入空间。

报告覆盖的三家核心标的中，GS对思源电气和国电南瑞给予买入评级，对华明装备维持中性。但比个股评级更重要的，是报告揭示的行业逻辑：在变压器交付周期长达128周的美国市场，思源电气能在6-9个月内完成交付，产能扩张周期仅需一年左右。这种速度优势，正在改写全球电力设备的竞争规则。

KC评论： 72\%的缺口听起来像是一个静态数字，但GS真正在说的是一个动态博弈——美国本土产能扩张需要2-3年，而中国供应商的扩产周期是一年。这中间的两年时间差，就是中国变压器企业最大的战略窗口。完整报告中对这个时间差的量化拆解，以及对各区域出口结构的详细分析，值得仔细推敲。

五月中国变压器出口增速从27\%跃升至72\%，表面看是全面开花，但拆解区域贡献后会发现，真正的结构性变化藏在结构里。

非洲贡献了252\%的同比增长，占出口总额的18\%；中东152\%的同比增长，占比23\%；美洲113\%的同比增长，占比21\%；亚洲35\%的同比增长，占比23\%；欧洲则是唯一的负增长区域，同比下降12\%，占比15\%。

这个结构告诉我们三件事。第一，中东和非洲的爆发式增长不是偶然，这两个区域的电力基础设施投资正在加速，且对价格敏感度较高，中国供应商的性价比优势在这里被放大。第二，欧洲市场的负增长值得警惕，可能反映了欧洲本土产能的恢复或贸易政策的调整。第三，美国市场虽然114倍的同比增长看起来很惊人，但这是在极低基数上的反弹，环比增长26\%，相比4月的95\%同比增速，实际动能并没有加速。

GS在这份报告中还提供了其他品类的出口数据：电子电表出口同比下降11\%，断路器出口同比下降16\%。这进一步印证了，变压器出口的强劲并非中国电力设备整体出口的普遍现象，而是特定品类的结构性机会。

KC评论： 很多人看到72\%的增速会认为中国变压器出口全面爆发，但GS的数据清楚显示，欧洲市场是负增长，电子电表和断路器也在下滑。这说明变压器出口的强劲更多是结构性需求驱动，而非中国制造业的整体竞争力提升。完整的区域出口拆分图表和品类对比数据，在报告中有更详细的呈现。

GS估算，美国电力变压器需求与本土供应的缺口将从72\%收窄至2028年的57\%。这个15个百分点的收窄，看似乐观，但背后是多重假设的叠加。

美国本土变压器制造商确实在扩产。Virginia Transformer刚刚宣布在阿拉巴马州新建一座变压器工厂。但问题在于，全球变压器交付周期仍然维持在约128周的高位，而新建产能从规划到投产通常需要2-3年。

相比之下，思源电气可以在6-9个月内交付变压器，产能扩张周期约一年。这种速度优势背后是中国供应链的敏捷性——上游原材料、零部件、熟练工人的配套能力，是美国本土制造商短期内难以复制的。

GS还指出，思源电气刚刚宣布了第三期变压器工厂，目标年产值达到90-100亿元人民币。这个规模对应的是全球化的产能布局思维，而不仅仅是满足国内需求。

KC评论： 128周交付周期与24-36周交付周期的差距，本质上不是生产效率的差异，而是整

[... middle omitted ...]

月以来一直稳定在高位，5月同比上涨3.9\%。而中国出口到美国的变压器价格，在10-220MVA规格上由于产品组合变化较大，难以观察出明确趋势。

这种价格分化说明，美国市场的定价权仍然掌握在本土供应商手中，中国供应商更多是在填补供需缺口，而非主动发起价格竞争。只要美国本土供需缺口维持在高位，中国供应商就有足够的议价空间。

原材料方面，变压器核心原材料取向硅钢（GOES）价格正在小幅下行，这对中国供应商的利润率是正向支撑。

KC评论： 29\%的价格跌幅很容易被解读为价格战信号，但GS的核心判断是“仍在正常区间”。这里的关键在于理解产品组合效应——不同规格、不同客户的变压器价格差异极大，单纯看均价变化可能会误判竞争格局。报告中对价格区间的详细图表和分规格分析，是判断价格趋势不可或缺的依据。

4. 报告尚未完全回答：中国变压器出口的关税和政策风险如何量化

这份GS报告提供了详实的数据和清晰的判断，但有一个关键问题没有完全展开：美国对中国变压器加征关税的可能性与影响。

\subsection{[R015] JPM：中国楼市真正的拐点，不在成交量，在挂牌量}
{\small\color{refgray}归类：地产与消费：中国楼市拐点在挂牌量，消费数据出现5月裂口}

这份JPM最新周度数据监测，在端午假期扰动之下，给出了一个值得认真对待的信号：中国房地产市场并非全面回暖，但一线城市二手房挂牌量的持续下降，正在为价格企稳构筑一个此前被市场低估的底部。这并非简单的“销售回暖”故事，而是一个关于供给收缩如何改变定价权的结构性变化。

市场普遍关注成交量，因为它是情绪的晴雨表。但JPM的数据显示，真正关键的变化发生在供给端。一线城市二手房挂牌量已从3月峰值下降2.6\%，上海单周挂牌量更是加速下降1.1\%。当卖盘不再涌出，买方才有机会在议价中不再被动。这正是本轮周期与以往最大的不同。

当然，端午假期打乱了所有周度数据的可比性。10城实时二手房成交同比从+13\%骤降至0\%，60城新房网签同比下滑23\%。JPM特别指出，若与去年同期端午假期对比，10城及一线城市成交分别上升15\%和10\%。数据噪音之下，趋势并未逆转。

但噪音本身也值得解读。假期效应暴露了市场的一个脆弱点：当前成交量的韧性高度依赖交易日的正常节奏，一旦遇到假期扰动，同比增速就会快速收窄。这意味着，市场尚未形成自我强化的内生动力。真正的考验在下半年，当政策脉冲效应消退，成交量能否在没有假期干扰的情况下维持正增长。

KC评论： JPM的数据提醒我们，不要被单一周的数据牵着走。端午假期造成的同比失真，在二手房网签数据中尤为明显——过去9周连续正增长的记录被中断。但剔除假期影响后，10城成交同比仍为正。这个细节在完整报告中有更细致的图表拆解，值得仔细对比。

JPM在信用债观点中明确提出“K型分化”正在加剧。5月统计局70城数据显示，一线城市二手房价格连续三个月环比上涨，但整体70城环比降幅从4月的-0.23\%扩大至-0.26\%。这组数据背后的含义非常清晰：市场正在分裂为两个世界。

一线城市的企稳，并非因为需求突然爆发，而是因为供给端出现了主动或被动的收缩。JPM跟踪的中原报价指数从17.8微降至17.3，说明业主涨价意愿并不强烈，但挂牌量下降意味着可售房源减少，买方选择范围收窄，议价空间自然压缩。

对于开发商而言，这意味着资产价值的重估将高度分化。在一线城市拥有优质土储和商业资产的公司，将受益于资产价格底部的确认。而广泛布局低能级城市的开发商，仍将面临持续的价格下行压力。JPM在股票和信用债层面的推荐，无一例外地指向了这种分化逻辑。

在JPM的数据中，上海的表现格外突出。一线城市二手房成交同比下滑3\%，但上海仍保持正增长。12城二手房网签同比下滑13\%，上海是唯一同比正增长的一线城市。更值得注意的是，上海二手房挂牌量单周下降1.1\%，降幅远高于其他一线城市。

这意味着上海正在形成“量缩价稳”的正向循环：挂牌量下降减少供给，成交韧性维持需求信心，价格预期逐步稳定。JPM的中原经理信心指数虽从53回落至51，但仍处于50荣枯线上方，表明市场参与者对后市并未转向悲观。

上海的特殊性在于其经济基本面和人口流入的支持。但这也提出了一个问题：其他一线城市能否复制上海的模式？北京、广州、深圳的挂牌量下降幅度明显小于上海，成交数据也更为波动。如果上海的经验无法被快速复制，那么全国市场的企稳节奏将比预期更为缓慢。

KC评论： 上海是当前中国楼市最值得观察的“实验室”。它的挂牌量下降速度领先于其他城市，这意味着价格修复可能也会领先。但报告没有展开的是，上海这轮挂牌量下降，有多少是政策限制的结果，有多少是市场自发出清？这个问题的答案，直接关系到其他城市的参考价值。完整报告中有更多城市层面的细分数据，可以帮助读者做出自己的判断。

JPM对香港住宅市场的判断非常直接：房价已实现全年10-15\%涨幅目标的上沿，预计下半年涨幅将放缓至5\%以下。这并非看空，而是对动能衰减的理性预判。

数据支撑了这一判断。香港房价指数周环比上涨0.6\%，年

[... middle omitted ...]

兆。

JPM的推荐组合也反映了这种谨慎：开发商方面，他们更青睐拥有优质资产和稳定派息能力的公司，如CKA和Sino；收租股方面，Swire Properties和HKL是首选。这暗示了他们对交易性机会的偏好正在减弱，而对现金流和资产质量的要求在提升。

4. 信用债市场的“K型”修复：新城控股的REITs操作值得关注

JPM信用债团队指出，JACI中国高收益地产指数上周上涨0.3\%，年初至今回报率达6.4\%。但真正值得关注的不是指数本身，而是个体信用之间的分化。

新城控股被单独提及。据Debtwire报道，新城控股计划在6月完成一笔20亿元人民币的私募REITs扩募，底层资产为天津和南通的两个购物中心。这笔操作如果成功，将覆盖该公司2026年剩余17亿元的在岸到期债务。JPM认为这一规模“可观”。

这揭示了一个重要的信用修复路径：拥有优质商业资产的公司，可以通过REITs等工具实现资产变现和债务覆盖，从而避免信用事件。而缺乏这类资产的公司，融资渠道将更为狭窄。

\subsection{[R016] JPM：存储芯片的“更高更久”周期，真正考验的是长协落地速度}
{\small\color{refgray}归类：AI/算力：需求三年翻五倍，MLCC与CoWoS成为新瓶颈}

JPM：存储芯片的“更高更久”周期，真正考验的是长协落地速度

过去三个月，亚洲存储芯片公司的股价上涨了44\%到184\%，同期费城半导体指数上涨20\%到88\%。这不是一个普通的周期波动。JPM在最新发布的存储市场研报中给出了一个明确的判断：存储芯片正处于一个“higher for longer”的上行周期，而这一轮周期的核心驱动力，已经从供需错配，转向了长期协议（LTA）的谈判进程与AI基础设施投资的再定价。

这份报告的价值，不在于它重申了“存储短缺”这个市场已知的事实，而在于它提供了一个观察框架：当所有人都看到短缺时，真正决定股价能否从“盈利上修”切换到“估值重估”的关键变量，是LTA何时落地。报告承认，LTA谈判比预期更慢，但预计2026年下半年将迎来密集宣布期。这正是当前市场情绪与基本面之间，最值得关注的裂缝。

*KC评论：** JPM的“higher for longer”不是一句口号。它的逻辑链是：短缺加剧 -> 长协锁定价格与量 -> 盈利可见性提升 -> 估值中枢上移。但当前市场只交易了前两步，第三步的“估值重估”还没有发生。这正是完整报告里最值得深挖的图表和假设。

5月以来，市场几乎没有听到新的存储长协公告。JPM认为，这并非需求出了问题，而是存储厂商在谈判条件上变得异常谨慎——价格条款、保护机制、预付款比例，每一项都在拉长谈判周期。三星的劳动罢工等公司治理事件，也分散了管理层的注意力。

但报告特别指出，美光与Anthropic的战略协议是一个值得关注的信号。这份协议不仅涉及HBM、DRAM和SSD的长期供应，还包括美光参与Anthropic的H轮融资。这种“股权投资+长期供应”的合同结构，与NAND SSD厂商和大型云服务商之间的LTA模式如出一辙。三星和SK海力士同样是Anthropic的战略股权伙伴，JPM预计，它们将在2026年下半年宣布与北美超大规模云服务商的类似协议。

这意味着什么？如果LTA在2026年下半年如期落地，存储芯片的定价模式将从“现货波动”转向“协议锁定”，这将直接支撑估值从当前的低位向历史中枢回归。报告认为，这才是存储板块真正的重估催化剂。

KC评论： 很多人把长协简单理解为“涨价”，但JPM的分析框架更精细：长协的本质是“将不可预测的现货波动，转化为可预测的盈利现金流”。一旦长协批量落地，存储公司的估值逻辑将从周期股切换为成长股。完整报告里对LTA合同结构的拆解，以及美光与Anthropic协议的细节分析，是理解这个判断的关键。

市场上一直存在一个担忧：随着AI模型效率提升（如DeepSeek、Turboquant），或网络架构优化（如CPO），单台服务器的存储容量需求可能会下降，从而缓解供需矛盾。

JPM的观点非常明确：当前存储内容优化的最大驱动力，不是模型效率提升，而是存储短缺本身。客户被迫接受12Hi HBM4E而非16Hi，不是因为不需要更高性能，而是因为16Hi根本买不到。SOCAMM（一种存储模组）的内容削减，更多是供应端瓶颈所致，而非需求端主动降级。

报告进一步指出，存储短缺没有缓解迹象，2027年可能比2026年更严峻。虽然部分投资者寄希望于网络侧的优化（如CPO规模化）来减少计算负载，从而降低存储需求，但JPM认为，当前的需求预测显示，几乎没有迹象表明需求破坏或内容优化会显著改善供需失衡。

*KC评论：** 这是报告中最反直觉的判断之一。市场普遍在讨论“DeepSeek会不会降低存储需求”，但JPM给出的答案是：短缺本身才是内容优化的原因，而短缺只会加剧，不会缓解。完整报告中对HBM4E 12Hi与16Hi的采样对比，以及对2027年供需缺口的量化测算，是支撑这个判断的核心证据。

3. 存储占CSP资本开支的比例，正在挑战AI

[... middle omitted ...]

云服务商愿意无限度地接受存储成本上升。

报告指出，过去12个月，存储价值在CSP资本开支中的占比从20\%快速攀升至50\%以上，主要发生在3月季度，当时DRAM/NAND价格的环比涨幅超出了年初预期。但JPM认为，半导体领域的盈利预测修正速度，远快于CSP领域的预测修正速度。两者之间的差距，可能需要几个月才能弥合。

KC评论： 这是整份报告最值得反复琢磨的图表（图1和图2）。存储价值占比从20\%到70\%的跳跃，本质上是在问：AI的商业模式，能否承受存储成为最大的单项成本？如果答案是“不能”，那么存储价格的上行空间将受到CSP采购意愿的制约。完整报告中对CSP资本开支与存储TAM的对比分析，以及历史周期的价值占比回归，能帮助你判断这个比例是否可持续。

报告特别提及，一家中国DRAM厂商已获得上海证券交易所的上市批准，计划融资295亿元人民币（约43亿美元），最早于7月上市。虽然具体细节尚未披露，但市场将密切关注其产能扩张计划、工艺良率和在传统DRAM市场的竞争策略。

\subsection{[R017] 摩根斯坦利：MS：2027年CoWoS的蛋糕，NVIDIA切走一半，但真正的变量在CPU}
{\small\color{refgray}归类：AI/算力：需求三年翻五倍，MLCC与CoWoS成为新瓶颈}

摩根斯坦利：MS：2027年CoWoS的蛋糕，NVIDIA切走一半，但真正的变量在CPU

当市场还在争论AI算力需求是否会在2026年见顶时，MS的一份最新报告给出了一个反直觉的判断：2027年，全球AI芯片的先进封装产能不仅不会过剩，反而会迎来新一轮结构性扩张，而推动这一轮增长的，不只有GPU。

这份由Charlie Chan领衔的亚太科技团队发布的报告，核心在于首次系统性地拆解了台积电2027年的CoWoS产能分配。结论清晰且有力：NVIDIA依然是最大的单一客户，但真正让这份报告值得细读的，是它揭示了CPU正在成为先进封装的新变量，以及Google TPU生态中联发科的意外崛起。

这不是一份简单的产能预测。它是一张2027年AI芯片竞争格局的路线图。谁在扩张，谁在收缩，谁在借势，都写在了CoWoS的订单里。

1. 2027年CoWoS总需求逼近270万片，NVIDIA占比降至45\%但绝对值仍翻倍

MS将台积电2027年底的CoWoS产能假设从17万片/月上调至20万片/月，这意味着全球AI XPU产业在2027年将实现约60\%的同比增长。而根据供应链核查，全球CoWoS总需求将从2026年的约139万片飙升至2027年的约269万片，接近翻倍。

NVIDIA的CoWoS消耗量预计增长57\%，达到122万片。虽然其份额从2026年的56\%下降到45\%，但这并非需求减弱，而是因为其他玩家的增速更快。NVIDIA在CoWoS-L上的消耗量增长40\%至91万片，主要服务于Rubin和Blackwell系列GPU。而CoWoS-R的订单激增则暗示了Vera CPU的出货量几乎翻倍——这一点与MS另一位分析师Joe Moore对NVIDIA数据中心收入增长52\%的预测一致。

KC评论： 份额下降不代表地位动摇。NVIDIA的绝对消耗量从78万片增至122万片，增量高达44万片，超过Broadcom或AMD的全年总量。这意味着，即便竞争对手在追赶，NVIDIA对先进封装产能的锁定能力依然是行业天花板。完整报告中有一张 年各客户CoWoS消耗量的堆叠图，直观展示了NVIDIA的领先优势如何被稀释但未被撼动。

2. Google TPU双供应商策略成型，联发科成为2027年最大变量

报告中最引人注目的变化来自Google TPU供应链。MS将联发科2027年的CoWoS-S预订量定为18万片，对应约360万颗TPU v8t芯片，远高于此前250万颗的出货假设。如果联发科能够帮助Google解决T-Glass基板供应问题，这一数字还有上行空间。

与此同时，Broadcom的CoWoS-S预订量为36.5万片，其中约33万片分配给TPU v8i，对应约390万颗芯片。这意味着，Google在2027年将同时通过联发科和Broadcom两条线生产TPU，且联发科的产品主打更高性价比的v8t，而Broadcom的v8i则拥有更大的芯片面积和更高的单颗价值。

更值得关注的是，报告披露了Google下一代2nm TPU的代号：TriggerFish。这是HumuFish的推理优化版本，暗示Google正在将训练和推理芯片分拆设计，以更精准地匹配不同工作负载。TriggerFish还可能用于TPU租赁服务，这将是Google云业务的一个潜在增长点。

KC评论： 联发科从手机芯片巨头跨界到AI ASIC，是这份报告中最有想象力的故事。18万片CoWoS-S的预订量意味着联发科在2027年将成为台积电CoWoS的第四大客户，仅次于NVIDIA、Broadcom和AMD。但这里有一个关键假设尚未验证：ABF基板的供应是否真的能跟上。报告中提到，联发科的出货量假设已经考虑了基板短缺的风险，但如果供应改善，上行空间可能超出

[... middle omitted ...]

027年的675万颗，增长超过5倍。Venice CPU的CoWoS预订量从5万片增至27万片，其中大部分产能由ASE/SPIL、Amkor和Powertech等非台积电阵营承接。

这一变化对供应链的影响是结构性的。AMD的CoWoS总消耗量增长308\%至53万片，其中非台积电部分占比从2026年的38\%提升至2027年的55\%。这意味着，ASE、Amkor等封测厂商正在从AMD的CPU扩张中获得增量订单，而不仅仅是等待NVIDIA的溢出。

KC评论： CPU进入先进封装，意味着AI算力的竞争已经从单纯的GPU军备竞赛，扩展到整个计算架构的重构。Agentic AI需要更频繁的推理和决策，而CPU在低延迟、高并发场景下的优势正在被重新定价。报告没有完全展开的是，这种CPU需求是否可持续，还是仅仅因为2027年是2nm CPU的换机大年。完整报告中有一张全球CoWoS产能按供应商拆分的图表，展示了台积电与非台积电阵营的产能扩张路径，是判断封测板块投资机会的核心素材。

\subsection{[R018] 摩根斯坦利：MS：稀土磁材的“中国之外”叙事，产能不是最大变量}
{\small\color{refgray}归类：稀土磁材与电力设备：中国之外产能规划漂亮但产出存疑，变压器交付速度成差异化优势}

摩根斯坦利：MS：稀土磁材的“中国之外”叙事，产能不是最大变量

当市场还在争论“中国稀土出口管制”会带来多大冲击时，MS近期与一位拥有35年行业经验的专家进行了深度对话，并发布了一份罕见的、聚焦于稀土永磁产业“中国之外”产能建设的研报。这份报告的核心判断，可能和许多人的直觉相反：美国规划的稀土永磁产能，到2028年理论上足以覆盖其当前的全部进口需求。但问题不在于能不能造出来，而在于造出来之后，谁用、用在哪、成本是多少。

这份报告的价值不在于给出了一个简单的“看好”或“看空”结论，而是系统性地拆解了从产能到成本、从技术到应用、从军用到民用之间的多层逻辑断层。对于关注产业链安全、新能源上游材料以及全球制造业格局重组的读者而言，这是一份必须细读的框架性文件。

1. 美国规划的产能数字很漂亮，但“产出”和“产能”是两回事

报告中最引人注目的数据是：美国钕铁硼（NdFeB）磁材产能到2028年可能达到约2.8万吨/年，如果加上已融资和潜在项目，这一数字可达3.5万吨，远期甚至可能超过5万吨。而美国当前从中国、越南和欧洲进口的散装磁材总量大约只有8500吨/年。单从数字看，美国规划的产能是其当前需求的3倍以上。

但专家在讨论中反复强调了一个核心区别：产能不等于产出。决定实际产出的，不是厂房和设备，而是客户认证、金属/合金供应、良品率、工艺控制和成本竞争力。对于有经验的企业，从工厂设计到满产大约需要3年；而对于初创公司，这个时间可能接近6年。这意味着，即便2028年的产能数字是真实的，其实际产出能否达到8500吨，仍然高度不确定。

KC评论： 这是一个经典的“规划泡沫”问题。投资者在看到大量产能建设公告时，需要区分“可融资的计划”和“可交付的产出”。完整报告中详细讨论了美国供应链中缺失的“氧化物到金属的转化”和“薄带铸轧合金生产”这两个关键环节，这些瓶颈是理解产能落地概率的关键。

2. 成本差距是西方磁材产业最硬的骨头，但质量差距并非不可逾越

专家给出了一个非常具体的判断：西方生产商大约需要3年时间就能达到中国磁材的质量水平，但在短期内几乎不可能达到中国的低成本。具体到价格，专家估计，西方生产的钕铁硼磁材可能比中国产品贵约50\%。

这个成本差距并非来自单一环节，而是贯穿于整个制造链条：从更低的人工成本、更成熟的供应链、到几十年的工艺学习积累。中国供应商在2000年代初期的定价就已经是美国售价的75\%，这直接导致当时美国生产商在盈亏平衡线上挣扎。一个具有竞争力的西方磁材产业，必须首先拥有一个能够支撑可持续定价的制造成本基础。

KC评论： 50\%的溢价在国防和某些高价值应用中是可以接受的，但对于价格敏感的电动汽车和风电行业，这是一个巨大的障碍。报告没有直接给出“成本差距何时缩小”的预测，而是提供了一套分析框架：需要关注哪些环节的成本下降路径。这正是完整报告中最值得深入推敲的部分。

市场经常讨论“减少稀土使用量”的趋势，但这份报告给出了一个精确的修正：过去二十年，钕铁硼磁材中的总稀土含量只从约32\%微降至约30.5\%。真正的进步发生在重稀土（镝和铽）的使用上。通过晶界扩散、晶界工程、更细的晶粒尺寸和更低的应用温度，某些应用场景中重稀土的需求已经减少了50\%以上。

而钕镨（NdPr）的替代则困难得多。专家明确指出，钕镨在钕铁硼磁材的化学结构中嵌入得更深。虽然镨的替代有助于拓宽可用原料来源，而铈/镧替代的牌号已经发展成一个低成本产品家族，但它们填补的是钕铁硼与铁氧体、铝镍钴等较低性能磁材之间的空白，而非替代高性能钕铁硼。

KC评论： 这一结论对投资者的含义是：不要因为看到了“稀土减量”的新闻就认为钕镨需求即将见顶。重稀土的减量故事是真实的，但钕镨的需求增长逻辑依然坚实。完整报告中对不同应用场景下磁材用量的具体

[... middle omitted ...]

航天/国防应用中仍然具有不可替代性。

KC评论： 500吨的国防需求与2万吨的规划产能之间形成了巨大的鸿沟。如果西方磁材产能只能依赖国防订单，那么产能利用率将非常低。真正的商业逻辑在于：除了国防之外，还有哪些民用领域愿意为“非中国”磁材支付50\%的溢价？报告没有直接回答这个问题，但它提供了一个分析框架：需要区分“必须去中国化”的应用和“可以继续使用中国产品”的应用。

报告中对终端市场的讨论非常务实。在美国，只有约1-2\%的已安装风力发电机使用钕铁硼磁材，欧洲的比例略高，而中国可能达到30-50\%甚至更高。在电动汽车领域，专家给出了一个具体数字：每千瓦额定电机功率的钕铁硼磁材重量约为12克，相当于每辆纯电动轻型车略高于2公斤，而混合动力车大约只用一半。

专家还提到，像移动机器人这样的新兴领域可能会增加需求，但许多应用使用的磁体非常小，虽然单价可能很高，但不需要数千吨的磁材产出。这意味着，对于大型磁材生产商而言，这些新兴市场可能无法提供足够的规模效应来消化新增产能。

\subsection{[R019] GS：美光财报暴露了一个被低估的欧洲AI供应链信号}
{\small\color{refgray}归类：AI/算力：需求三年翻五倍，MLCC与CoWoS成为新瓶颈}

当一家存储芯片公司季度收入同比增长346\%，并给出超出市场预期15\%的指引时，市场通常只看到了“存储周期向上”这一个故事。但GS欧洲半导体团队在这份最新研报中，揭示了一个更值得关注的信号：美光财报背后，欧洲AI赋能者正在进入一个需求可见性显著提升的阶段。

美光3QFY26收入约415亿美元，同比增长346\%，环比增长74\%，远超市场预期的363亿美元。更关键的是，公司预计4QFY26收入将达到500亿美元，环比再增21\%。数据中心收入单季超过250亿美元，折合年化运行率约1000亿美元。这些数字本身已经足够震撼，但GS分析师Alexander Duval及其团队真正想告诉市场的是：这些数字对欧洲半导体设备商和AI基础设施提供商意味着什么，而不仅仅是存储芯片本身。

这份报告的核心判断是：美光的业绩和指引，为欧洲半导体设备股ASML、ASMI、BESI以及AI基础设施提供商NBIS提供了需求可见性改善的明确信号，但不同公司的受益逻辑和程度存在本质差异。

1. 存储供需紧张从周期性波动转向结构性锁定，设备商的订单可见性正在质变

美光在财报中明确表示，AI正在推动数据中心的强劲增长，内存需求显著高于供应。公司预计这种供需失衡将持续到2027年之后，供应要到2028年才逐步改善。这不是一个典型的存储周期顶部信号，而是一个AI驱动的结构性短缺。

GS团队关注到一个关键细节：美光已经签署了16份“照付不议”战略客户协议，覆盖数据中心、消费和汽车终端市场。这些协议不仅包含价格下限，还包含价格上限，五年内累计承诺收入在保底价格下约为1000亿美元。更值得注意的是，这些合同覆盖了美光预计DRAM出货量的约20\%和NAND出货量的三分之一。

KC评论： 存储行业过去十年最痛苦的教训就是“周期波动”。但照付不议协议加上价格上下限，意味着美光未来五年的收入能见度发生了质变。对设备商而言，这直接转化为客户更愿意分享技术路线图和生产节奏的细节。设备商不再需要猜测客户何时扩产，而是可以基于可见的合同需求来规划交付。这是整个欧洲半导体设备板块估值逻辑的一个潜在升级点。

GS团队指出，存储供需紧张对欧洲半导体设备股ASML、ASMI和BESI构成短期至中期利好，因为供给侧紧张使得客户更愿意分享技术和产品爬坡的细节。换句话说，设备商的订单可见性正在从“季度博弈”转向“年度规划”。

2. ASML是欧洲半导体设备股中受益最直接、最确定的标的，原因在于其存储敞口和长交付周期

GS明确表示，ASML在其覆盖的欧洲半导体设备股中，处于最有利的位置来受益于这一需求背景。原因有两个：相对较高的存储芯片敞口和超过12个月的EUV工具交付周期。

ASML约40\%的收入来自存储芯片客户，这是欧洲半导体设备商中存储敞口最高的。更关键的是，美光最近与ASML签署了多年期EUV供应协议，支持其在1-delta节点及未来世代增加EUV采用。这与GS团队对EUV层数在未来几年将持续增加的判断一致。

长交付周期在这里意味着什么？当客户看到未来数年的需求确定性上升时，他们必须提前下单锁定产能。ASML的EUV工具交付周期超过12个月，这意味着今天的订单锁定将转化为未来一年以上的收入可见性。这与那些交付周期较短的设备商形成鲜明对比。

KC评论： GS给ASML的目标价是1770欧元，基于40倍CY27市盈率。这个估值倍数并不便宜，但关键在于这个倍数能否维持。如果存储客户的照付不议合同模式扩展到设备采购，ASML的收入可见性将从“季度波动”升级为“年度可预测”，估值中枢可能上移。完整报告中包含了ASML目标价的历史调整路径和估值框架的详细拆解，这是理解当前估值是否合理的关键材料。

3. BESI是HBM需求持续走强的隐藏受益者，混合键合技术正在从概念验证走向

[... middle omitted ...]

未来数年混合键合需求的关键组成部分，持续强劲的需求将有利于这一技术的采用。BESI的客户支出周期性和混合键合采用延迟是主要风险，但当前信号偏向积极。

4. NBIS的AI基础设施故事正在从概念走向执行验证，电力容量扩张是关键跟踪指标

GS团队将NBIS（Nebius）纳入这份报告的交叉验证范围。美光关于AI需求强劲的评论对NBIS同样构成利好，因为NBIS拥有裸金属和软件产品组合，并在扩展合同电力容量方面取得了令人鼓舞的进展。

NBIS作为欧洲AI基础设施提供商，其核心逻辑是：企业级AI工作负载需要从公有云向专用基础设施迁移，而NBIS提供的就是这个中间层。美光的数据中心收入年化运行率超过1000亿美元，只是整个AI基础设施需求的一个缩影。NBIS的合同电力容量扩张速度，是验证其需求是否真实、是否可持续的关键指标。

GS给NBIS的目标价是267美元，基于9倍CY27E EV/Sales倍数。这个倍数在AI基础设施公司中并不极端，但关键取决于收入增长能否兑现。

\subsection{[R020] GS：日本造船业被低估的三个催化剂，不只是AI的替代品}
{\small\color{refgray}归类：供应链与运价：高运价从事件冲击转向供给约束，中国变压器交付速度成优势 / 区域市场：人民币中间价模型释放偏强信号，日本造船与热泵迎来结构性机会}

当全球资金疯狂涌入AI主题时，一个传统行业正在东京湾和长崎的船坞里悄然积累着未来三年的利润。GS在最新发布的日本造船业深度报告中，给出了一个与市场主流情绪截然相反的判断：日本造船股的盈利前景实际上正在向上移动，而市场对AI相关资产的追逐，反而制造了估值洼地。

这份报告覆盖了名村造船、三井E\&S和东京计器三家GS覆盖标的，同时追踪了日本发动机、古野电气、中国涂料等多家非覆盖公司。核心结论是：日本造船业已经锁定了未来三年半的订单，而即将到来的政策催化剂——产能扩张与防务需求——可能将这一景气周期从“周期性修复”升级为“结构性重估”。

市场目前对日本造船股的谨慎态度，主要源于资金流向AI主题的挤出效应。但GS认为，这种谨慎恰恰忽略了基本面正在发生的实质性改善。截至2026年5月底，日本造船业的订单积压量达到2927万总吨，相当于未来三年半的产能已被填满。2026年前五个月的累计订单量同比仅微降2\%，仍处于历史高位。

KC评论： 市场习惯把造船看作周期性行业，但GS这份报告的核心逻辑是，日本造船业正在经历从“周期”到“结构”的转变。订单积压不是问题，问题是产能不够——而这恰好是接下来政策要解决的事情。完整报告里对产能扩张的具体路线图、各公司受益程度的量化拆解，才是真正值得细读的部分。

日本国土交通省海事局在2026年2月底宣布修订造船相关政策，启动产能扩张框架。关键条款是：企业必须制定从2024年水平扩张50\%以上产能的计划，才有资格获得造船产业复兴基金的支持——该基金总额3500亿日元，为期10年。日本内阁预计在2026年夏季前后批准包含造船在内的17个战略产业增长战略。

这意味着什么？目前日本造船厂已经满负荷运转。如果产能扩张计划落地，订单量和订单积压将不仅反映当前产能，而是包含未来可扩张的产能空间。这种预期一旦形成，将成为提升盈利增长前景和估值的催化剂。

GS特别指出，市场关注的焦点将在于名村造船等主要船厂是否会利用复兴基金。这不仅是产能问题，更是行业信心的试金石。

日本防务需求的催化剂密集排列在2026年下半年：6月的经济财政运营基本方针、7月的防卫白皮书、年底前修订三项战略安保文件、以及美国2027财年国防预算。

其中最具想象空间的是美国国防预算。媒体报道称，美国可能拨出约18.5亿美元用于利用日本或韩国船厂采购海军舰艇。如果订单落在日本，GS覆盖标的东京计器作为舰艇部件供应商将直接受益。同时，名村造船的船舶维修业务也将从日美海军舰艇维修需求增长中获益。

KC评论： 防务需求对日本造船业来说不是“锦上添花”，而是“第二曲线”。名村造船的维修业务、东京计器的部件供应，这些都不是一次性订单，而是伴随海军舰队寿命周期的持续收入流。完整报告里对防务需求在各公司收入中的占比、以及对利润率的影响有更细致的拆解，这部分数据对于判断估值弹性至关重要。

全球造船业订单趋势方面，GS注意到伊朗局势导致运输距离拉长、替代采购增加，运费飙升，原油油轮需求尤其强劲。新船价格仍处于历史高位。

这是一个典型的“地缘政治→运费→新船订单”传导链条。运费上涨意味着船东现金流改善，进而有能力接受更高的新船价格。而日本造船厂目前订单积压已满，议价能力处于强势地位。

但GS也提醒，这种短期需求驱动因素需要与产能扩张、防务需求等中长期催化剂配合，才能支撑持续的盈利增长。

GS追踪了多家非覆盖公司的最新动向，这些信息对于判断整个产业链的景气度至关重要：

日本发动机预计在3/27财年实现稳定的销售和利润增长，驱动力是中国发动机制造商从2026年第三季度开始的产能扩张带来的授权费收入增长。古野电气向欧洲航运公司和中国船厂的销售强劲，防务业务也确认了利润率改善的潜力。中国涂料则面临成本压力——年初以来氧化亚铜和环氧树脂价格上涨，主

[... middle omitted ...]

所有环节都能同步受益。完整报告里对各公司业务结构和利润率驱动因素的对比分析，能帮助读者判断哪些环节的议价能力更强。

5. 报告尚未完全回答的关键问题：产能扩张的节奏和利润率弹性

GS这份报告虽然给出了明确的看多判断，但有几个关键问题仍需进一步验证。

第一，产能扩张的具体节奏。日本造船厂目前满负荷运转，扩张产能需要时间、资本和劳动力。政策框架虽已推出，但企业实际扩产计划何时落地、产能爬坡速度如何，报告没有给出明确时间表。

第二，利润率弹性在不同公司之间的差异。名村造船受益于大型船舶的利润率改善，三井E\&S的发动机业务受益于国内供应紧张，东京计器受益于防务需求。但各公司的利润率改善幅度、可持续性、以及产能扩张对资本开支和折旧的影响，需要更细致的量化分析。

第三，中国造船厂的竞争动态。GS报告中提到中国造船厂产能扩张对日本发动机厂商的授权费收入是利好，但中国造船厂本身也在全球订单竞争中占据份额。日本造船厂能否在产能扩张的同时维持价格纪律，是一个需要持续观察的问题。

\subsection{[R021] GS：热泵在英国卖的不是取暖，是空调}
{\small\color{refgray}归类：区域市场：人民币中间价模型释放偏强信号，日本造船与热泵迎来结构性机会}

这份研报的核心判断，是在一个出奇炎热的英国夏天里形成的。当GS的分析师团队在伯明翰的InstallerSHOW上与欧洲热泵协会（EHPA）负责人以及五家核心设备供应商逐一交流时，他们捕捉到了一个正在发生的结构性转变：全球变暖正在推动空气-空气热泵（air-to-air heat pump）在英国和欧洲的加速渗透。这不是一个关于冬季取暖的故事，而是一个关于夏季制冷需求爆发的故事。对于长期关注欧洲能源转型的投资者而言，这意味着热泵市场的竞争逻辑、产品形态和估值框架，都需要因此重新校准。

这份报告的价值，不仅在于它提供了英国热泵市场的最新一手数据——2025年销量约12.5万台，同比增长27\%，渗透率仍不足1\%——更在于它揭示了几个正在同时发生、且彼此强化的趋势：政策补贴的加码、中国供应商的激烈竞争、以及欧洲本土OEM通过产品线扩张和渠道护城河来应对的策略。以下是我们从这份研报中提炼出的五个关键洞察。

1. 全球变暖正在改写热泵的产品逻辑：制冷功能成为新的增长引擎

GS在研报中明确指出，全球变暖很可能会推动空气-空气热泵在英国和欧洲的进一步普及。这个判断的现场感极强——报告提到，在InstallerSHOW举办的那一周，英国和欧洲正经历异常高温。在这种天气下，热泵的制冷功能不再是锦上添花的附加项，而是变成了刚需。

传统上，英国中央供暖系统围绕“湿回路”（即水循环散热片）设计，空气-空气热泵在2025年英国市场的销量占比不到10\%。但EHPA总干事Paul Kenny向GS团队透露，这种情况可能会随时间改变，因为整个欧洲的空气-空气热泵市场规模大约是空气-水热泵的6倍。更关键的是，多个OEM正在积极响应安装商的反馈和规模经济需求，大幅拓宽产品线。NIBE最近宣布进入其历史上从未涉足的空气-空气热泵市场，Viessmann也在InstallerSHOW上宣布在英国推出其Vitocal 200-A型号。

KC评论： GS的这个判断，实际上在提醒投资者：不要再用“取暖设备”的老框架来给热泵公司估值。如果制冷功能成为标配，热泵的市场空间将从“替代锅炉”升级为“替代空调+锅炉”，目标市场的天花板将显著抬升。但这背后有一个关键问题尚未回答：空气-空气热泵在英国现有建筑中的改造安装成本，是否真的能降到消费者可接受的水平？完整报告中包含了EHPA对各国安装成本差异的详细拆解，这部分数据对于判断市场渗透速度至关重要。

2. 英国热泵市场的政策窗口正在收窄，但补贴结构的变化是更精准的信号

英国热泵市场的长期增长潜力毋庸置疑，但GS报告中最值得关注的，是政策补贴结构的变化。英国政府的“温暖家园计划”目标是到2030年实现每年约45万台的安装量。目前的核心工具是锅炉升级计划（BUS），为空气源热泵或地源热泵安装提供7500英镑补贴。从今年7月21日起，使用石油或液化石油气的离网住户，安装热泵还可额外获得1500英镑补贴。

但真正值得注意的，是两个结构性变化。第一，补贴范围已扩大到包括排气空气热泵，这对NIBE是明确的利好，因为它是英国市场上仅有的三家供应商之一。第二，从今年晚些时候开始，空气-空气热泵也将获得2500英镑的补贴资格。这等于政策层面正式承认了制冷功能的价值。

不过，GS也提醒，英国热泵市场在2026年第一季度表现疲弱，原因是低收入资助计划转型带来的过渡期扰动。NIBE预计2026年英国市场整体销量持平。这意味着短期政策红利并不均匀，不同产品线和不同定位的公司，受益程度将显著分化。

3. 中国供应商的竞争不是“会不会来”，而是“怎么打”——渠道成为欧洲OEM的核心护城河

报告对竞争格局的讨论非常具体。GS团队观察到，中国HVAC OME和其他新进入者正试图通过低成本模式或直销模式来获取市场份额，长期来看可能对

[... middle omitted ...]

果只是靠价格优势做D2C，可能很难撼动欧洲OEM的渠道壁垒。但报告也提到，市场正在出现“增量向D2C模式移动”的趋势，包括英国天然气公司与NIBE的合作。渠道之争的胜负，最终取决于谁能更好地服务安装商——而安装商驱动了超过90\%的住宅购买决策。完整报告中包含了IMI如何通过其Heatmiser品牌和超过150万个联网设备来建立“拉动式需求”，这部分内容对理解渠道壁垒的数字化升级非常有价值。

EHPA总干事Paul Kenny向GS团队透露了一个关键数字：要保持竞争力，热泵OEM的年产量需要超过50万台。这个门槛对大多数欧洲中小型OEM来说，是巨大的挑战。

报告指出，热泵OEM的数量远多于生产关键零部件的公司。例如，全球只有四家主要供应商生产压缩机，而压缩机占热泵总价值的20\%-25\%。这种供应链结构意味着，下游OEM的整合将是必然趋势。Kenny预计，未来几年市场将进一步整合。亚洲和北美OEM对欧洲市场的渗透野心，以及规模经济对产量的硬性要求，将共同推动这一进程。

\subsection{[R022] NOM：字节跳动的音乐App，正在用AI绕过版权战争}
{\small\color{refgray}归类：未被主题摘要引用}

字节跳动旗下音乐平台Soda Music的DAU已逼近5000万，年增长仍高达50\%，但它走的是一条与腾讯音乐截然不同的路：不烧钱买独家版权，靠广告变现，甚至默许AI翻唱绕过付费墙。NOM最新专家调研揭示，这家公司真正的野心不在音乐本身，而在成为AI音乐时代的规则制定者。

这份报告最值得关注的判断，并非Soda Music的用户增长数字，而是其背后字节跳动对“音乐”这一品类的战略定位：一个用来消耗过剩广告库存、补齐内容生态短板的辅助型护城河，而非独立盈利中心。这与市场普遍将其视为“腾讯音乐挑战者”的叙事形成了根本差异。

*NOM分析师在报告中提出一个关键问题：当版权预算受限时，AI音乐能否成为Soda Music的弯道超车点？** 答案或许比表面看起来更复杂。

KC评论： 这份报告的价值不在于告诉你Soda Music有多少用户，而在于揭示了字节跳动内部两个音乐App（Soda和番茄音乐）在相互竞争，且用户重叠度高达60-70\%。这意味着字节跳动在音乐赛道的真实投入，可能远小于外界想象。完整报告中还有关于两个App用户画像的详细拆解，以及字节跳动对合并计划的明确表态。

1. 用户增长正在从“流量输血”转向“活动拉新”，边际效益递减已是定局

截至2026年5月，Soda Music的DAU约为5000万，MAU为1.5亿。专家预计年底DAU将达到约6000万，同比增长约50\%。这看似亮眼，但与2025年近100\%的增速相比，减速信号已十分明显。

*核心原因在于：抖音的流量转化池正在接近饱和。** 那些可以被抖音自然导流的“可转化用户”基本已被洗过一遍。剩余的用户要么缺乏使用音乐App的刚性需求，要么已经被其他娱乐形式（如小说、短剧）截胡。

Soda Music的应对策略是转向线下：赞助音乐节、校园活动，瞄准低线市场。这些活动由赞助商买单，成本可控，但拉新效率显然无法与抖音的免费流量相比。

用户日均使用时长也从高峰期的60-70分钟下降至约50分钟。NOM专家指出，这既是用户基数扩大带来的稀释效应，也反映了来自其他在线娱乐活动的竞争压力。

*这一趋势对字节跳动的含义是：音乐App的用户增长天花板已经清晰可见，未来增长必须依靠产品本身而非渠道红利。**

2. 广告收入占比70\%，Soda Music本质上是一个“广告播放器”而非音乐订阅服务

这是NOM报告中最反常识的发现之一。Soda Music的收入结构中，广告占比约70\%，订阅收入仅约30\%。这与腾讯音乐形成鲜明对比——后者在线音乐收入的60-70\%来自订阅。

*这种结构并非偶然，而是字节跳动有意为之。** 专家明确表示，字节跳动对Soda Music的核心要求是“维持盈亏平衡”，而非不惜代价追求增长。Soda Music的角色被定义为抖音生态系统中的“辅助护城河”——满足用户的听歌需求，同时消化抖音剩余的广告库存。

目前，Soda Music在剔除人员成本和总部成本后已达到运营盈亏平衡，但全口径计算仍处于亏损状态。

KC评论： 这意味着任何关于Soda Music“颠覆腾讯音乐”的叙事都需要重新审视。它根本就不是奔着取代QQ音乐去的。它的存在逻辑更像是字节跳动版的“用户留存工具”——就像微信的“看一看”功能，不是为了打败今日头条，而是为了不让用户离开微信。完整报告中还有Soda Music与腾讯音乐在版权覆盖率上的具体对比数据，以及专家对版权格局未来走势的判断。

Soda Music的曲库覆盖了全国消费歌曲的约75\%，低于腾讯音乐的90-95\%和网易云音乐的85\%。在这75\%中，60\%是授权版权音乐，15\%是平台自制和UGC内容。

*版权预算的严格约束，决定了Soda Music不可能在“独家版权战争”中与腾讯正面交锋。** 但

[... middle omitted ...]

的用户欢迎。对于这些侵权内容，Soda Music仅在收到权利方完整证据的删除通知后才进行处理，没有主动审查义务。

这形成了一种“猫鼠游戏”：权利方的执法资源有限，无法处理所有侵权内容。这种模糊策略，实际上降低了Soda Music的内容获取成本，同时保持了对低端用户的吸引力。

4. 粉丝经济与线下演唱会：Soda Music明确选择“不参与”

NOM专家给出了一个清晰的战略取舍：Soda Music不会涉足大型偶像演唱会、K-pop粉丝活动等资源密集型业务。

原因有三： 监管层面对粉丝文化的压力持续收紧 利润率极薄 缺乏能将抖音流量转化为可货币化粉丝互动的专业运营团队

取而代之的是，Soda Music将继续执行小规模、由赞助商资助的线下活动，如三四线城市的音乐节和校园活动。这些活动的定位是“成本可控的用户获取渠道”，而非深度参与粉丝经济。

*这一选择进一步强化了Soda Music的“辅助工具”定位：它不愿意为一个低利润、高风险的业务投入核心资源。**

\subsection{[R023] 摩根斯坦利：MS：中国通胀的真相不在PPI，而在工资条}
{\small\color{refgray}归类：地产与消费：中国楼市拐点在挂牌量，消费数据出现5月裂口}

这可能是当下最反直觉的一个判断：中国近期转正的PPI，并非通胀回归的信号。

MS在最新发布的亚洲经济报告中，直接点破了这个市场普遍存在的误解。报告明确指出，最近几个月PPI同比转正，主要驱动力是油价上涨，而非国内需求的实质性复苏。随着油价回落，这个“伪通胀”信号将很快消失。

真正值得关注的，不是宏观价格指数，而是微观层面的三个指标：非大宗商品部门的利润率、工资增速、以及零售销售。这三个指标，才是判断中国能否从“通缩”走向“温和再通胀”的钥匙。

这份由MS首席亚洲经济学家Chetan Ahya和首席中国经济学家邢自强领衔的报告，系统拆解了中国再通胀进程中的周期性与结构性力量。核心结论是：出口驱动的周期性改善正在发生，但结构性障碍仍未解除，一个2-3\%的常态化通胀还不在视野之内。

市场对通胀的讨论，往往被PPI等价格指数牵着走。但MS提醒投资者，必须区分“价格回升”和“需求驱动的通胀”。

报告指出，PPI在经历了约三年半的通缩后转正，但这并非可持续通胀的信号。过去几个月PPI的上行动力主要来自油价上涨。随着国际油价近期回落，这个动力正在消退。更关键的是，非大宗商品部门的利润率并未改善——换句话说，企业仍然没有重新获得定价权。

与PPI形成鲜明对比的是零售销售数据。报告特别提到，随着消费品以旧换新政策效果消退，5月零售销售已经出现绝对收缩。酒店和交通领域的数据同样疲软：6月酒店每间可用客房收入（RevPAR）同比仅增长2\%，尽管去年同期基数极低；端午节期间的航空客运增长也相当乏力。

KC评论： 零售数据是理解中国当前经济状态最直接的窗口。PPI可能被上游价格扭曲，但零售数据直接反映消费者的真实意愿。当零售销售在低基数下仍然收缩，说明问题不在供给端，而在需求端。完整报告中提供了详细的图表，展示了零售、酒店、航空等多个高频指标的同步走弱，值得仔细对照。

MS观察到，中国政策制定者正在重复一个熟悉的剧本：当外部需求强劲时，他们会逆周期地收紧政策刺激力度。

今年一季度实际GDP增长超预期，叠加出口持续强劲，政策层在二季度明显放缓了财政支出节奏。数据显示，广义财政赤字占GDP的比例从1月的11.9\%收窄至5月的10.9\%。反映到固定资产投资上，名义FAI增速降至年内低点-10.7\%，其中基础设施FAI从一季度的9.2\%骤降至5月的-10.8\%。

这种操作模式并非首次出现。报告指出，在过去几轮周期中，每当外需强劲，政策层倾向于减少内需刺激，以避免经济过热和通胀压力。但这一次，二季度GDP增速已回落至4.4\%，低于全年增长目标。MS中国经济学团队预计，政策层将在三季度重新加快财政支出节奏，以平滑增长。

KC评论： 这种“外需好就收、外需弱就放”的操作模式，本质上是一种对国内需求的不自信。它意味着政策层更倾向于用出口作为增长引擎，而非主动刺激内需。完整报告详细拆解了这种模式的历史规律和背后的决策逻辑，对于理解未来政策节奏至关重要。

3. 结构性障碍比周期性压力更棘手：人口、地产与社保三重约束

如果说周期性问题可以通过政策调整来缓解，那么结构性问题则需要更根本性的改革。MS指出了三个相互交织的结构性因素，它们共同压制了中国的通胀潜力。

第一，人口结构恶化。中国15-64岁劳动年龄人口已从2015年峰值的10亿降至目前的9.9亿，总人口自2022年起持续下降。这构成了对总需求的长期结构性拖累。

第二，房地产市场调整。弱化的人口趋势意味着房地产需求的结构性下降。叠加政策支持的力度有限，房价持续承压。但报告也指出，房地产调整已经相当深入：房地产固定资产投资（不含土地成本）占GDP的比重已从2014年峰值的16\%降至2025年的6\%，甚至低于2004年水平。这意味着房地产对经济的拖累将逐步减弱，但也不太可能成为未

[... middle omitted ...]

杆，但也是最难推进的。因为增加社会福利支出意味着永久性的财政承诺，而当前政府收入占GDP比重持续下降，财政空间有限。完整报告对这三个因素的相互作用机制有更深入的分析，特别是关于“为什么社保改革迟迟难以推进”的讨论，值得仔细阅读。

既然结构性改革短期内难以期待，那么推动再通胀的周期性力量就变得更加重要。MS认为，亚洲的工业和资本支出超级周期正在驱动出口复苏，而中国是主要受益者之一。

报告指出，亚洲出口复苏正在从科技领域扩展到非科技领域，驱动因素包括能源、国防、AI及相关基础设施、工业供应链建设等多个领域的资本支出。中国的非半导体出口在过去两个月刚刚恢复到接近两位数的增长。

这种出口改善将如何传导到国内通胀？逻辑链条是：出口增长提升产能利用率，进而改善非大宗商品部门的利润率，然后企业才有信心和能力提高工资，工资增长最终支撑消费。但报告也提醒，这个传导存在滞后——工业产能利用率刚刚从一季度低位开始改善，去产能进展仍然缓慢，因此出口改善传导到利润率和工资增长还需要时间。

\subsection{[R024] NOM：人民币中间价模型已释放一个明确信号}
{\small\color{refgray}归类：区域市场：人民币中间价模型释放偏强信号，日本造船与热泵迎来结构性机会}

人民币兑美元中间价的定价模型，正在发生值得关注的偏移。NOM亚洲外汇策略团队最新发布的模型测算显示，其核心预测值较前次下调了152个基点，计入逆周期因子后的预测值也较前次下调了124个基点。这不是一次简单的技术调整，而是模型层面透露出一个信号：在外部压力与内部目标之间，政策天平正在微妙移动。

这份报告的价值不在于给出一个具体的汇率预测数字，而在于它拆解了影响中间价的两个关键变量——模型本身如何定价，以及逆周期因子如何介入。对于关注人民币资产定价、跨境资本流动和出口企业汇兑风险的读者而言，理解这两个变量的变化方向，比知道6.80还是6.81更重要。

NOM模型的最新预测值为6.8043，较前次官方预测值6.8195低了152个基点。更值得注意的是，这并非一次性的技术修正。报告附带的模型误差图显示，近期模型预测值与实际中间价之间的偏差已经持续了一段时间。误差不是随机波动，而是呈现出方向性。

这意味着什么？模型是根据一系列可量化的输入变量——包括隔夜美元指数变动、主要货币对走势、市场情绪指标等——来生成预测的。当模型持续高估或低估实际中间价时，通常有两种解释：一是模型遗漏了某些关键变量，二是政策当局在模型之外施加了额外的影响。NOM的这份报告通过同时呈现“不含逆周期因子的模型预测”和“含逆周期因子的模型预测”，实际上为读者提供了一个观察窗口：哪些变化来自市场，哪些变化来自政策。

KC评论： 理解这份报告的关键不在于记住6.8043这个数字，而在于看懂它比前次预测低了152个基点。这意味着模型认为人民币应该比此前预期的更强。如果读者只看最终中间价数字，会错过模型层面方向性变化的信号。完整报告中的误差图可以直观展示这种趋势是否在加速。

NOM报告提供了两组预测值：不含逆周期因子的模型预测为6.8043，含逆周期因子的预测为6.8071。两者之间的差值仅为28个基点。这个差值就是当前逆周期因子对中间价的“加码”幅度。

对比历史数据，逆周期因子在人民币贬值压力较大时期，其调节幅度可以高达数百个基点。当前28个基点的差值，说明政策当局对中间价的主动干预已经降到极低水平。这不是因为外部压力消失，而是因为政策逻辑发生了变化——在当前的汇率形成机制下，模型本身已经在朝着政策合意的方向运行，逆周期因子不再需要大幅介入。

这种变化对市场参与者有直接含义：中间价的波动将更多地反映模型输入变量的变化，而非政策的临时调整。这意味着隔夜美元走势、一篮子货币变动等外部因素对次日中间价的影响权重在上升。对于做市商和外汇交易员而言，模型预测的参考价值在提升，政策不确定性的溢价在下降。

3. 2026年下半年的事件日历构成了人民币汇率的关键压力测试

NOM报告附带的未来事件日历，表面上只是一份时间表，实际上是一份“压力测试清单”。7月政治局会议、10月国庆黄金周、11月深圳APEC会议、12月中央经济工作会议和政治局会议，以及年底前中国国家主席习近平对美国的访问——这些事件密集排列在下半年，每一件都可能成为影响汇率预期的催化剂。

尤其值得关注的是年底前的中美元首会晤。报告引用特朗普在2月8日采访中的表态，称习近平主席将在年底前访问白宫。这一事件的特殊性在于，它既可能成为缓解贸易紧张关系的契机，也可能因为双方预期落差而加剧市场波动。在APEC深圳会议之后紧接着安排元首访问，时间窗口非常紧凑。

对于企业财务管理者而言，这意味着下半年的汇率风险管理窗口正在收窄。如果等到事件结果明朗再调整对冲头寸，可能已经错过了最佳时机。报告没有直接给出交易建议，但事件日历本身已经暗示了波动率的潜在上升路径。

KC评论： NOM这份报告最被低估的价值在于它提供了事件的时间轴。很多市场参与者只关注模型数字，忽视了这些事件对汇率预期的重塑作用。完整报告中还

[... middle omitted ...]

型预测值可能指向人民币贬值方向。但与此同时，国内经济复苏需要稳定的汇率预期来吸引外资流入，APEC会议和元首访问也需要营造积极的合作氛围。在这种情况下，是让模型继续主导定价，还是重新加大逆周期因子的介入力度？

NOM报告提供了观察框架，但没有给出答案。这恰恰是读者需要自己判断的部分。模型的预测能力在正常市场环境下是可靠的，但在政策目标与市场信号出现冲突时，模型的边界就会显现。理解这个边界在哪里，比理解模型本身更重要。

基于NOM这份报告的分析逻辑，读者在跟踪人民币汇率走势时，应该建立以下三个观察维度，而不是仅仅关注每日中间价的绝对水平。

第一，模型误差的方向和幅度。如果模型持续高估或低估实际中间价，且误差在扩大，说明有模型之外的力量在起作用。这时需要判断是逆周期因子的介入，还是模型本身需要调整。

第二，逆周期因子的变化趋势。当前28个基点的差值处于历史低位，但如果这个差值开始扩大，意味着政策当局在重新校准干预力度。扩大的方向（是抑制贬值还是抑制升值）同样重要。

\subsection{[R025] DB：中国新能源车订单的“冰火两重天”已到关键分水岭}
{\small\color{refgray}归类：新能源车与电池：订单分化加速，电池行业进入三国分治阶段}

中国新能源汽车市场正在经历一个罕见的“订单分裂”时刻。DB最新发布的周度订单追踪数据显示，2026年6月第三周，头部车企的周订单量呈现出截然不同的方向：比亚迪单周订单突破10万辆，环比暴增134\%，而蔚来、华为鸿蒙智行（AITO）等品牌却出现两位数下滑。这不是简单的月度波动，而是竞争格局加速分化的信号。

这份报告的核心判断是：在价格战常态化、消费者预期趋于理性的背景下，订单数据的“冰火两重天”正在揭示一个更深层的结构性变化——规模效应和品牌心智的差距，正从利润表传导至订单簿。那些能持续将流量转化为订单、将订单转化为交付、将交付转化为口碑的车企，正在获得指数级的优势。而这一趋势，可能在未来2-3个季度内决定中国新能源车行业的最终座次。

为什么现在关注这个信号？因为6月第三周的数据恰好处于年中冲刺的关键节点。车企的促销政策、新车型上市节奏、以及消费者对补贴退坡的预期，都在这个时间窗口集中释放。DB的周度订单追踪，是市场上少数能实时捕捉这些动态的高频指标。它比月度销量数据更早揭示趋势，比市场情绪更接近真相。

1. 比亚迪单周破10万不是意外，而是规模效应的“临界点”爆发

6月第三周，比亚迪新订单达到10.68万辆，环比增长134\%，同比增长34\%。这是一个令人瞩目的数字。但更值得关注的是环比增速：前一周（6月第二周）订单仅为4.57万辆，而去年同期为8万辆。这种“前低后高”的节奏，反映出比亚迪在年中促销节点上的精准发力。

从结构上看，比亚迪的订单暴增并非来自单一车型的爆款效应，而是其“王朝+海洋”双网体系的全面开花。这种规模效应一旦形成，就会产生正向循环：更多的订单意味着更强的供应链议价权，更强的议价权意味着更低的成本，更低的成本又支撑更激进的价格策略。DB的数据显示，比亚迪的订单波动性正在降低，这意味着其客户基础已经从“价格敏感型”转向“品牌忠诚型”。

KC评论： 比亚迪的订单数据说明，在新能源车行业，规模已经不是护城河，而是“入场券”。真正的护城河在于，当规模达到某个临界点后，能否将规模转化为对上游供应商的议价权和对下游消费者的定价权。比亚迪正在证明它可以。但其他车企能否复制？答案藏在它们各自的订单结构里。

2. 蔚来订单同比暴增195\%，但环比骤降26\%——“以价换量”的可持续性存疑

蔚来6月第三周订单为1.45万辆，同比暴增195\%，但环比下降26\%。这个反差值得深究。195\%的同比增速确实亮眼，但需要注意到去年同期基数极低（仅0.49万辆）。环比26\%的下滑，则暗示促销活动的边际效应正在递减。

蔚来的问题在于，它的订单增长高度依赖新车型和促销政策的短期刺激。一旦促销力度减弱或竞品推出更有吸引力的方案，订单就会迅速回落。这种“脉冲式”增长模式，在行业整体增速放缓的背景下，很难转化为可持续的竞争优势。

更关键的是，蔚来的订单波动率在所有追踪车企中是最高的之一。这意味着其客户群对价格和促销的敏感度极高，品牌溢价尚未真正建立。DB的图表显示，蔚来的周订单曲线在5月下旬至6月初经历了一轮快速拉升，但随后迅速回落，呈现出典型的“脉冲”特征。

3. 理想汽车订单同比降34\%，环比仅微增8\%——中高端市场的“天花板”正在显现

理想汽车6月第三周订单为0.59万辆，同比下降34\%，环比仅增长8\%。这是所有追踪车企中同比降幅最大的。理想的问题不在于产品力，而在于其目标市场——30万元以上中大型SUV——正在遭遇需求瓶颈。

在价格战向中高端市场蔓延的背景下，理想的增程式技术路线虽然解决了里程焦虑，但未能有效抵御来自纯电车型和混动竞品的挤压。华为问界（AITO）的崛起，直接分流了理想的核心客户群。数据显示，问界6月第三周订单为1.17万辆，虽然环比下降36\%，但同比仍增长30\%，远优于理想。

KC评论：

[... middle omitted ...]

企中仅次于比亚迪。小米的优势在于其强大的品牌号召力和线上线下融合的销售网络。但需要警惕的是，小米的订单基数仍然较低，且其增长主要来自SU7一款车型。

小米的“后发优势”体现在：它不需要像蔚来、小鹏那样从零开始建立品牌认知，小米的粉丝群体可以直接转化为潜在购车用户。但“后发劣势”同样明显：它进入市场时，消费者已经对新能源车有了更成熟的认知和更高的期望值。小米能否持续推出有竞争力的新车型，将是决定其订单增长可持续性的关键。

DB的报告没有直接给出结论，但图表显示，小米的订单曲线在5月中旬经历了一轮调整后，6月第三周再次回升。这种“W型”走势，暗示其订单增长可能受到产能和交付节奏的制约。

5. 特斯拉订单同比仅增5\%，环比降4\%——全球霸主在中国市场的“舒适区”正在缩小

特斯拉6月第三周订单为1.26万辆，同比增长5\%，环比下降4\%。这个增速在所有追踪车企中是最低的之一。特斯拉在中国市场的增长放缓，并非产品力下降，而是因为其“降价-订单-再降价”的策略循环正在失效。

\subsection{[R026] 摩根斯坦利：中国消费的“5月裂口”比表面数据更值得警惕}
{\small\color{refgray}归类：地产与消费：中国楼市拐点在挂牌量，消费数据出现5月裂口}

*一份来自摩根斯坦利的五月消费追踪报告，揭示了两个趋势：消费活动指数进一步下滑，而零售总额同比增速自2022年以来首次转负。但真正值得关注的，不是这些数字本身，而是它们共同指向的一个判断——支撑消费的短期政策效果正在衰减，而结构性压力尚未见底。**

对于关注中国资产配置的投资者来说，这意味着什么？如果你仍在期待一场“消费复苏”来支撑企业盈利的V型反弹，这份报告提供了足够多的理由需要重新审视这个假设。

摩根斯坦利的中国消费活动Z-Score，是一个基于五个高频指标的复合指数，覆盖了居民贷款、餐饮零售、商品零售（不含汽车）、乘用车销售和航空客运量。这个指数在五月继续走低，而MSCI中国指数的同比变化也同步下行。这不是一个孤立的月度波动——它是连续趋势的延续。

KC评论： 摩根斯坦利的Z-Score并非简单的零售数据加总，它的价值在于捕捉了“消费意愿”和“消费能力”的同步变化。当航空客运、餐饮和商品零售同时走软时，这说明消费者不仅在减少大额支出，连日常可选消费也在收缩。这才是真正值得警惕的信号。

*零售总额同比增速在五月转为-0.6\%，这是自2022年以来的首次负增长。** 报告指出，造成这一变化的核心原因有两个：以旧换新政策的刺激效果正在消退，以及就业市场的持续疲软。换句话说，政策驱动的消费脉冲已经释放完毕，而内生增长动力尚未接棒。

对于企业决策者而言，这意味着一个更严峻的竞争环境正在形成。当整体蛋糕不再增长甚至缩小时，市场份额的争夺将变得更加残酷。那些依赖消费升级逻辑的商业模式，可能需要重新审视自己的定价权和增长路径。

*报告对未来给出了明确的判断：消费可能进一步放缓。** 支撑这一判断的三个理由是：政策支持力度减弱、劳动力市场条件疲软、以及出口繁荣对消费的溢出效应有限。

这里有一个容易被忽视的细节。市场上有一种观点认为，出口的强劲表现会通过企业利润和就业传导至消费端。但摩根斯坦利的报告暗示，这种溢出效应可能被高估了。出口部门的复苏主要集中在特定行业和头部企业，并未广泛惠及整体就业市场。这意味着，消费的修复不能依赖出口的“间接拉动”，而需要更直接的内生动力。

KC评论： 这份报告最有价值的地方不在于它说了什么，而在于它没有说什么。它没有给出“触底”的时间表，没有提出“政策必须出手”的建议，也没有预测“下半年会好转”。它只是冷静地展示了数据，并给出了一个基于逻辑的判断。这种克制本身就是一种态度——在不确定性面前，承认不确定性比强行给出结论更有价值。

*那么，报告没有完全回答的关键问题是什么？** 至少有两个。

第一，政策工具箱中还有哪些选项可以扭转消费颓势？报告提到了“政策支持减弱”，但没有深入分析如果消费进一步恶化，决策层是否会推出新的刺激措施，以及这些措施可能的形式和效果。对于投资者来说，这恰恰是最需要跟踪的变量。

第二，消费结构内部是否存在分化？报告使用的是总量指标，但不同收入阶层、不同城市层级、不同品类的消费表现可能存在显著差异。如果高端消费和大众消费出现分化，那么对资产配置的含义是完全不同的。高端消费品的投资逻辑和大众消费品的投资逻辑，在当前的宏观环境下需要分开评估。

*对于读者来说，这份报告提供了一个观察框架：** 与其关注月度数据的波动，不如关注三个核心指标——就业市场的边际变化、政策刺激的力度和节奏、以及出口繁荣是否能够真正传导至消费端。当这三个信号同时出现改善时，才是真正可以重新评估消费板块的时机。

报告的完整版本包含了更多细分数据，包括每个因子的具体走势、历史对比、以及摩根斯坦利对中国消费板块的具体配置建议。这些细节对于构建投资假设至关重要。例如，Z-Score的每个构成因子都在报告中有独立的图表和解读，可以帮助读者判断哪些领域是真正拖累消费的“短板”，哪些领域出现了边

[... middle omitted ...]

商品零售都在减速。零售总额同比-0.6\%，2022年以来首次转负。

2️⃣ 两个微弱亮点 居民贷款小幅回暖 乘用车零售略有改善 但这两个改善的幅度，远不足以对冲整体下行。

3️⃣ 背后原因 以旧换新政策效果在消退，就业市场持续疲软，出口对消费的拉动也有限。研报判断，后续消费可能继续放缓。

💡 几个观察角度 消费降级不是短期现象，结构性压力在积累 政策刺激的脉冲效应越来越短，市场需要更实质的支撑 出口好≠消费好，内外需传导链条在断裂

China Equity Strategy | Asia Pacific

China Five-factor Consumer Activity Z-Score vs. MSCI China

\textbackslash{}- Consumer Z-score declined further in May as air passenger traffic, catering, and commodity retail sales all softened.

\subsection{[R027] DB：中国奢侈品进口的“五月寒”比数据本身更值得警惕}
{\small\color{refgray}归类：地产与消费：中国楼市拐点在挂牌量，消费数据出现5月裂口}

2026年二季度开局，中国奢侈品进口数据给了市场一记清醒的耳光。

DB最新发布的《China luxury import: Slower April-May vs. months of 1Q》报告显示，2026年4月至5月，中国奢侈品进口量同比下滑7.2\%，而一季度这个数字还是正增长5.0\%。这不是一个简单的季节性波动。进口量的环比恶化发生在市场普遍预期“消费温和复苏”的背景下，恰恰说明驱动奢侈品消费的核心力量——财富效应和消费者信心——正在重新走弱。

这份报告最值得关注的判断不是“奢侈品卖不动了”，而是：修复进程中的中国奢侈品市场，正在被两个同时出现的负向力量拖回原点——房价持续承压叠加股市财富效应快速消退。

一季度进口数据的改善，曾被部分投资者解读为消费信心的底部修复。DB的数据表明，那更像是春节错位和低基数的短暂共振。二季度头两个月的数据，正在把叙事拉回到一个更现实的框架里。

KC评论： 一季度+5.0\%和二季度前两个月-7.2\%的落差，不是统计噪声。它意味着那些基于“中国消费V型复苏”假设做多奢侈品股票的仓位，需要重新审视其核心逻辑。完整报告中有按品类、按国别拆分的详细进口量变化图表，能帮你判断哪些品牌的暴露最危险。

1. 手袋进口从+5.1\%骤降至-11.2\%，是二季度最弱的品类

从品类结构看，手袋是二季度最受伤的板块。4-5月手袋进口量同比下降11.2\%，而一季度还是正增长5.1\%。其中，非皮革手袋从一季度的+7.5\%降至-12.4\%，皮革手袋从+2.5\%降至-9.8\%，跌幅几乎对称。

更值得玩味的是5月单月数据。皮革手袋进口量同比下滑20.6\%，是过去一年中最差的一个月。这个数字比4月的+0.9\%急转直下，说明问题不是“缓慢恶化”，而是“突然失速”。

从产地看，意大利手袋进口量二季度同比下滑9.0\%，法国手袋下滑14.0\%。意大利手袋通常对应Gucci等品牌，法国手袋对应爱马仕、LV等。两者同步走弱，说明这不是某一个品牌或某一类消费者的选择问题，而是整体需求收缩。

KC评论： 进口量数据是品牌在中国市场真实需求的“硬指标”，因为它不受渠道库存和定价策略的干扰。DB报告中有手袋进口量与LVMH、爱马仕、Gucci亚太区同店销售的对比图，历史相关性极高。如果你持有这些公司的股票，那几张图值得仔细看。

珠宝品类呈现一个有趣的分化：进口量同比-8.6\%（一季度为-6.0\%），跌幅扩大但幅度有限；而零售端金银珠宝销售额从一季度的+10.6\%急转直下至-17.1\%。进口和零售的巨大背离，意味着渠道库存正在快速累积，终端需求比进口数据显示的更弱。

从月度趋势看，珠宝进口量在5月为-7.4\%，略好于4月的-9.7\%，但3个月移动平均仍为-15\%。DB报告指出，珠宝进口数据与相关品牌亚太区同店销售具有较好的季度相关性，这个信号需要认真对待。

手表是唯一在5月出现改善的品类：瑞士手表对亚洲出口量5月同比+2.9\%，而4月为-6.9\%。但二季度整体仍为-1.9\%，远低于一季度的+15.9\%。这个改善是否可持续，需要6月数据确认。

3. 真正令人担忧的不是进口数据本身，而是宏观背景的同步走弱

DB报告的核心贡献，不在于提供了最新的进口数据，而在于把进口数据放进了财富效应和消费者信心的框架里。

5月全国房价同比下跌4.8\%，与4月的4.9\%基本持平。房价没有继续恶化，但也没有改善。更重要的是，股市财富效应正在快速消退：上证综指和恒生指数5月同比涨幅从4月的21.0\%收窄至14.8\%。对高净值消费者而言，股票账户的浮盈缩水，会直接抑制大额可选消费的意愿。

消费者信心指数（最新数据为4月）同比仅增长1.4\%，而一季度平均为3.3\%。这个指标在3月已经放缓至2.9\%，4月进一步走弱。DB的结论很直

[... middle omitted ...]

.0\%，3月+3.4\%。2月的强劲很大程度上是春节错位和2025年低基数的结果。3月已经开始放缓，只是被一季度的整体数字掩盖了。

零售端的数据更清晰：一季度服装鞋帽零售+6.5\%，金银珠宝+10.6\%，看起来消费在恢复。但4月这两个数字分别降至+2.0\%和-23.1\%，5月进一步降至+1.1\%和-8.9\%。一季度的高增长更像是一次性的补偿性消费，而非趋势性改善。

DB报告没有直接说“中国奢侈品消费没有复苏”，但数据已经给出了这个判断。

5. 报告尚未完全回答的关键问题：是“暂时性调整”还是“趋势性逆转”？

支持“暂时性调整”的论据：二季度前两个月的数据可能受到4月清明假期和5月劳动节假期出境游增加的影响——消费者在海外购买奢侈品，导致进口数据偏低。如果这个假设成立，三季度进口数据可能反弹。

支持“趋势性逆转”的论据：房价和股市财富效应同时走弱，消费者信心持续下行，这些结构性因素不是假期效应可以解释的。如果三季度宏观数据没有明显改善，奢侈品进口可能进一步恶化。

\subsection{[R028] GS：奶粉行业最坏的时候正在过去，但赢家已非旧面孔}
{\small\color{refgray}归类：地产与消费：中国楼市拐点在挂牌量，消费数据出现5月裂口}

中国奶粉行业正在经历一场“量价双杀”的尾部出清，但GS最新追踪的尼尔森与线上数据揭示了一个更微妙的信号：行业总量的收缩幅度在边际上趋于稳定，而真正的分化发生在品牌之间——不是国产与进口的简单二分，而是渠道能力、库存周期与品牌信任度的重新排序。

这份发布于2026年5-6月的研报，覆盖了截至2026年4月的线下数据与5月的线上数据。对于关注中国消费品的投资者而言，核心判断只有一个：行业底部可能正在形成，但下一轮增长的受益者，不会是上一轮周期的全部赢家。

GS引用的尼尔森数据显示，2026年3-4月，中国婴幼儿配方奶粉线下市场同比下降8.2\%，较1-2月的7.7\%降幅有所扩大。但深入拆解后，信号并不完全悲观。

关键在于结构：3-4月线下销量同比下降9\%，较1-2月的8.7\%降幅仅扩大0.3个百分点；而平均售价同比微增0.9\%，较1-2月的1.1\%增幅略有收窄。这意味着，行业的下行主要仍由销量驱动，而非价格战加剧。母婴店渠道下降7.3\%，现代商超渠道下降20\%，渠道分化持续。

KC评论： 量跌价稳的组合，说明行业正在经历“被动缩量”而非“主动降价”。这比量价齐跌更接近底部信号。但完整报告中的图表显示，线上渠道5月降幅扩大至19\%，京东渠道更是同比下滑26\%，这提示线上库存压力可能仍在释放中。想判断线上是否先于线下触底，需要看完整报告中的季度增长趋势图。

伊利线下渠道在3-4月实现1.9\%的同比增长，较1-2月的-0.9\%显著改善，且跑赢行业-8.2\%的降幅。驱动因素来自母婴店渠道2.8\%的增长。伊利在母婴渠道的份额持续提升，说明其渠道下沉与品牌信任度正在转化为实际销售。

飞鹤的情况则更为复杂。线下渠道3-4月同比下降12\%，虽然较1-2月的-15\%有所收窄，但仍显著跑输行业。飞鹤的线下市场份额从1-2月的20.6\%微升至3-4月的21\%，同比仍下降0.9个百分点。线上渠道更为严峻，5月天猫/淘宝/京东合计同比下降31\%，尽管较4月的-47\%有所改善，但基数效应不可忽视。

KC评论： 伊利转正与飞鹤仍负增长的对比，是这份报告最值得追问的地方。伊利做对了什么？是产品结构升级，还是渠道管理更精细？飞鹤的份额回升是否可持续？完整报告中对两家公司的估值与风险部分有更详细的拆解，包括对出生率、竞争格局、政策支持等关键变量的敏感性分析。

GS的数据显示，3-4月线下市场外资品牌整体价值份额从1-2月的23.5\%回升至23.8\%，增加0.3个百分点。但外资内部的分化同样剧烈。

A2 Milk在母婴店渠道实现2.9\%的同比增长，市场份额从4.1\%升至4.3\%，同比增加0.4个百分点，主要靠销量增长4\%驱动。达能（包括多美滋与纽迪希亚）线下份额小幅提升。但惠氏（雀巢）线下同比下降60\%，份额从1-2月的1.3\%降至1.1\%，同比降幅达1.5个百分点。

线上渠道同样呈现“强者恒强”的格局。5月天猫/淘宝/京东合计，爱他美（达能）同比增长6\%，市场份额维持在20\%左右；贝拉米（蒙牛）同比增长19\%，尽管增速从4月的80\%回落；而A2线上同比下降25\%，可能与其4月起主动控制供应有关。

KC评论： 外资品牌整体份额回升，但惠氏和纽迪希亚的大幅下滑说明，消费者对外资品牌的“信任溢价”正在收窄，只有真正具备产品差异化或渠道优势的品牌才能守住阵地。完整报告中的线上市场份额变化图表，可以帮助判断哪些品牌正在“吃”掉竞争对手的份额。

GS的线上数据揭示了一个值得警惕的现象：5月行业线上销售额同比下降19\%，但京东渠道同比降幅高达26\%。考虑到京东是奶粉线上销售的重要渠道，这一降幅可能不仅仅是需求疲软，更可能反映了渠道库存的去化加速。

研报提到，5月线上降幅扩大部分源于去年同期的高基数（618大促）。但更深层的问题是：如果线上库

[... middle omitted ...]



第二，渠道结构是分化的第一层过滤器。伊利在母婴渠道的持续增长，飞鹤在线上渠道的剧烈波动，都说明不同品牌在不同渠道的竞争壁垒差异巨大。单纯看“全渠道”数据可能会掩盖关键信号。

第三，ASP趋势是判断竞争烈度的核心指标。如果ASP开始普遍下滑，意味着行业进入价格战阶段，这将压缩所有玩家的利润空间。目前ASP仍为正增长，但线上折扣力度需要持续跟踪。

第四，库存周期是短期波动的放大器。无论是线上还是线下，渠道库存的去化速度都会影响短期销售数据。GS的数据显示，部分品牌的线上波动可能与主动控制供应有关（如A2），而非需求崩溃。

KC评论： 这些观察框架在完整报告中都有对应的图表和数据支撑，包括按渠道、按品牌、按月度的详细追踪。对于希望建立自己判断框架的读者，完整报告中的图表比文字更有价值。

飞鹤的线上降幅收窄是趋势反转还是基数效应？伊利能否维持母婴渠道的增长势头？外资品牌中，谁最有可能复制A2的份额提升路径？线上库存去化何时完成，以及它对线下ASP的传导效应有多大？

\subsection{[R029] DB：德国养老金改革的最大变量不是延迟退休，而是强制入市}
{\small\color{refgray}归类：宏观与利率：ECB加息未止，德国养老金改革重塑长期资金流}

德国养老金改革专家组近日发布了33项改革提案。大多数分析聚焦于“延迟退休至67.5岁”或“取消提前退休优惠”——这些确实是重要信号，但并非最值得关注的变量。

这份提案中真正具有结构性意义、且可能对全球资产配置产生深远影响的，是一个被多数中文媒体低估的条款：引入强制性资本积累制，将2个百分点的毛工资收入强制投入资本市场，预计每年带来350亿欧元的增量资金。

这不是一个“养老金调整”的新闻，而是一个关于“德国是否会成为下一个瑞典式资本强国”的底层逻辑变化。

DB在最新发布的研报中，将这一条款称为“game-changer”——改变游戏规则的元素。而这一判断，可能比你想象的更值得认真对待。

DB研报明确指出，瑞典的“premium pension”体系是德国此次改革的直接模版。瑞典模式的核心逻辑是：在现收现付制（PAYG）之外，强制要求每个劳动者将一部分收入投入资本市场，通过长期复利积累个人养老金资产。

德国提案的具体设计是：从2028年起，每年逐步增加0.5个百分点的强制缴存比例，到2031年达到2个百分点。这笔资金由雇主和雇员各承担一半，投入一个集中管理的公共养老金基金。

关键数字：到2031年，每年新增资本市场的投资额将达到300至350亿欧元。

KC评论： 350亿欧元是什么概念？它相当于德国DAX指数成分股公司年分红总额的一半左右。而且，这还只是“强制部分”的规模。DB估计，同期改革后的私人养老金（需主动加入）可能带来每年约500亿欧元的资金流入。两者合计，德国每年将有超过800亿欧元的新增资金涌入资本市场。对于一个长期被视为“低风险偏好、高储蓄率、低股票配置”的经济体，这一转变的边际影响不可小觑。

更值得关注的是，这笔资金的资产配置策略尚未确定。DB研报指出，养老金委员会仅提到“投资于资本市场”，没有明确股票与固定收益的占比。但报告同时强调，要实现足够的长期回报以提升养老金水平，“必须高度聚焦股票市场”——这与瑞典模式的经验一致。

这就意味着，未来几年，德国可能从一个“全球股票配置偏低”的典型代表，转变为“机构化长期资金持续流入”的新玩家。对于全球资产管理行业，这可能是下一个十年最重要的结构性增量之一。

多数媒体对德国养老金改革的报道集中在“退休年龄从67岁延至67.5岁”——这确实是一个明确的信号，表明德国政府正在正视老龄化带来的财政压力。

但DB研报的分析框架显示，真正区分“修补”与“重构”的，是制度层面的变化。

现收现付制（PAYG）的稳定措施包括：重新引入“可持续因子”（自动减少养老金涨幅当缴费者减少时）、取消“63岁免减额退休”等。这些措施的作用是“稳定”而非“增长”——它们延缓了PAYG系统崩溃的时间，但没有创造新的价值来源。

而强制资本积累制的引入，则意味着德国养老金体系从“单一支柱”开始向“多支柱”转型。这是一个根本性的制度重构。

KC评论： 这就像一家公司从“只靠主营业务现金流发工资”转向“同时建立投资部门，用资本收益补充员工福利”。前者是成本控制，后者是资产创造。DB研报的核心洞察在于：德国政府终于意识到，仅靠提高缴费率或延迟退休无法解决老龄化社会的根本矛盾——劳动人口减少意味着现收现付制的税基萎缩。唯一的出路是让资本积累成为第二引擎。

DB研报坦诚地指出了几个关键的不确定点——这些不是报告的缺陷，而是政策制定过程中尚未落地的变量。对于希望从中推导投资逻辑的读者来说，这些不确定点恰恰是最值得跟踪的信号。

*第一，公共基金的管理架构。** 德国目前没有一个主权财富基金或公共资产管理机构。研报提到，建立这样一个“基础设施”将是关键——可能依托德国央行、KfW（德国复兴信贷银行）或Kenfo（德国养老金管理机构）。但具体由谁管理、治理结构如何设计、投资决策是否

[... middle omitted ...]

当前PAYG系统的缴费率已从18.6\%预计升至2040年的21.8\%。新增的2个百分点强制资本积累是“额外缴费”——这意味着德国劳动者的总社保负担将从18.6\%升至约23.8\%。这笔新增成本是否会影响消费、储蓄行为，以及是否会导致资本外流或劳动力市场扭曲，报告没有展开，但这是每个关注德国经济的人都需要思考的问题。

长期以来，全球投资者对德国的印象是：一个依赖出口、人口老龄化、数字化落后的经济体。德国股票市场在全球配置中的权重与其经济规模并不匹配。

如果改革按计划推进，德国将出现一个持续、可预测、规模巨大的长期资金流入资本市场。这对于德国本土的资产管理行业、交易所、以及那些能够提供稳定分红和长期增长的企业，都是一个结构性利好。

同时，这也意味着德国政府正在从“财政纪律优先”转向“资本市场建设优先”——后者是英美模式的核心特征之一。如果这一转变持续，德国可能在未来十年经历一轮“居民资产从银行存款向股票基金迁移”的过程，类似于日本在2014年NISA改革后出现的情况。

\section{Disclaimer}
{\scriptsize\color{refgray}
This community is a paid learning and information-sharing community created for general educational purposes only. The membership fee is charged solely for access to community discussions, educational content, organization, and operational costs. It is not an advisory fee, management fee, performance fee, brokerage fee, or compensation for personalized financial, investment, legal, tax, or professional advice.\par
I participate and share information solely as an individual and for educational discussion among community members. I am not acting as a financial adviser, investment adviser, broker, dealer, portfolio manager, legal adviser, tax adviser, accountant, fiduciary, or any other licensed professional adviser. Nothing shared in this community constitutes, and should not be interpreted as, financial advice, investment advice, legal advice, tax advice, a recommendation, solicitation, offer, endorsement, instruction, or guarantee to buy, sell, hold, short, trade, or invest in any security, cryptocurrency, fund, derivative, strategy, or other financial product.\par
All content shared in this community, including messages, discussions, links, files, charts, examples, market commentary, watchlists, AI-generated or AI-polished summaries, and third-party materials, is general, impersonal, and educational in nature. It does not take into account any individual's financial situation, investment objectives, risk tolerance, investment horizon, liquidity needs, tax status, legal restrictions, or suitability requirements. No content should be relied upon as suitable for any specific person, account, portfolio, or financial decision.\par
Information shared in this community may be incomplete, inaccurate, outdated, speculative, AI-generated, AI-edited, or based on third-party internet sources that have not been independently verified. I make no representation or warranty as to the accuracy, completeness, timeliness, reliability, or suitability of any information shared.\par
Financial markets involve substantial risk, including the possible loss of principal. Past performance, historical data, hypothetical examples, simulated results, backtests, personal opinions, or educational examples do not guarantee future results. Each member is solely responsible for conducting their own research, exercising independent judgment, and making their own decisions. Before making any financial, investment, legal, or tax decision, members should consult qualified and licensed professionals.\par
Participation in this community, payment of any membership fee, or communication with me or other members does not create any adviser-client, fiduciary, professional, contractual, agency, partnership, employment, or other advisory relationship. By joining or participating in this community, you acknowledge and agree that you are solely responsible for your own decisions, actions, transactions, profits, losses, risks, and consequences, and that I shall not be liable for any direct, indirect, incidental, consequential, financial, legal, tax, or other loss or damage arising from or related to any content, discussion, information, or material shared in this community.\par
}
\end{document}
